%% toc: P3 36 36.1,36.3
\section{Proofs: Mostowski quantifiers and the uncountability quantifier}
\label{sec:annex-p03-w03-p03-042}

Notation is that of \cref{sec:p03-w03-p03-042-mostowski}.  Where the
axiom of choice is used it is named.  ``Computable'' and ``recursively
enumerable'' for sets of formulas are read through an effective numbering
of expressions, and a described decision or listing procedure is
identified with a recursive one by the Church--Turing thesis
(\cref{def:p04-church-turing}).

\begin{annexitem}{annex:p03-w03-p03-042-support}{thm:p03-w03-p03-042-support}{Supports exist}
\begin{proof}
Let \(x,y\in I^{*}\) agree on \(K_1\cap K_2\).  Define \(z\in I^{*}\) by
\(z_j=x_j\) for \(j\in K_1\) and \(z_j=y_j\) otherwise.  Then \(x\) and
\(z\) agree on \(K_1\), so \(F(x)=F(z)\).  For \(j\in K_2\): if
\(j\in K_1\) then \(j\in K_1\cap K_2\) and \(z_j=x_j=y_j\); if \(j\notin
K_1\) then \(z_j=y_j\).  So \(z\) and \(y\) agree on \(K_2\) and
\(F(z)=F(y)\).  Hence \(F(x)=F(y)\).  The family of finite sets with the
property is nonempty by definition and closed under pairwise intersection;
a member of least cardinality is included in every member (otherwise its
intersection with that member would be smaller), so it is the least one.
\end{proof}
\end{annexitem}

\begin{annexitem}{annex:p03-w03-p03-042-theorem-one}{thm:p03-w03-p03-042-theorem-one}{Mostowski's Theorem~1}
\begin{proof}
\emph{\(Q_T\) is a quantifier.}  Let \(\varphi\) be a permutation of
\(I\) and \(F\) a one-argument function, identified with
\(X=F^{-1}(\vee)\subseteq I\).  Then \(F_\varphi\) is identified with
\(\varphi[X]\) and \(F_\varphi^{-1}(\wedge)=\varphi[I-X]\).  Since
\(\varphi\) is a bijection, \(\overline{\varphi[X]}=\overline X\) and
\(\overline{\varphi[I-X]}=\overline{I-X}\), so \(Q_T(F_\varphi)=Q_T(F)\).
Also \(\overline X+\overline{I-X}=\overline I\), so the pair belongs to
the domain of \(T\).

\emph{Existence of \(T\).}  Let \(Q\) be limited to \(I\) and
\((\mathfrak m,\mathfrak n)\) a pair with \(\mathfrak m+\mathfrak n=
\overline I\).  By the definition of the cardinal sum there are disjoint
sets \(A,B\) with \(\overline A=\mathfrak m\), \(\overline B=\mathfrak n\)
and a bijection \(f\) of \(A\cup B\) onto \(I\).  Put \(X=f[A]\); then
\(I-X=f[B]\), so subsets \(X\) with \(\overline X=\mathfrak m\) and
\(\overline{I-X}=\mathfrak n\) exist.  If \(X,Y\) are two such subsets,
choose bijections \(g:X\to Y\) and \(h:I-X\to I-Y\); \(\varphi=g\cup h\) is
a permutation of \(I\) with \(\varphi[X]=Y\), so by invariance \(Q\) takes
the same value at \(X\) and at \(Y\).  Define \(T(\mathfrak m,\mathfrak
n)\) as this common value.  Then \(Q_T(F)=T(\overline X,\overline{I-X})=
Q(F)\) for every \(F\).

\emph{Uniqueness.}  If \(Q_T=Q_{T'}\), then for each pair
\((\mathfrak m,\mathfrak n)\) take \(X\) as above:
\(T(\mathfrak m,\mathfrak n)=Q_T(X)=Q_{T'}(X)=T'(\mathfrak m,\mathfrak n)\).

\emph{Dual.}  For \(X\subseteq I\),
\(Q_{T^*}(X)=T^{*}(\overline X,\overline{I-X})={\sim}T(\overline{I-X},
\overline X)={\sim}Q_T(I-X)\), and \(I-X\) is the set of \({\sim}F\).
\end{proof}
\end{annexitem}

\begin{annexitem}{annex:p03-w03-p03-042-denumerable-classification}{thm:p03-w03-p03-042-denumerable-classification}{Quantifiers on a denumerable set (Mostowski, \S2(e))}
\begin{proof}
\emph{Item 1.}  If \(\overline I=\aleph_0\) and
\(\mathfrak m+\mathfrak n=\aleph_0\), then not both are finite, so the pair
is \((n,\aleph_0)\), \((\aleph_0,n)\) with \(n\) finite, or
\((\aleph_0,\aleph_0)\).  By Theorem~1, \(Q_T\) is determined by \(T\),
hence by these values.

\emph{Item 2.}  For \(Q_T\) limited to \(I\), the sequence
\(T(n,\aleph_0)\) is constant from some \(n\) on iff \(A_T\) is finite or
cofinite in \(N\), that is, iff \(A_T\in\mathfrak A\); for \(\mathfrak A\)
generated by \(\mathfrak I\) consists exactly of the finite and the
cofinite subsets of \(N\), these being closed under complement and finite
union and containing the finite sets.  Likewise for \(B_T\).  So the map
\(h(Q_T)=\langle A_T,B_T,\gamma_T\rangle\) sends \(\mathcal A\) into
\(\mathfrak C\).  It is injective by item~1.  It is onto: given
\(\langle A,B,\gamma\rangle\in\mathfrak C\), define \(T(n,\aleph_0)=\vee\)
iff \(n\in A\), \(T(\aleph_0,n)=\vee\) iff \(n\in B\),
\(T(\aleph_0,\aleph_0)=\gamma\); then \(Q_T\in\mathcal A\) and
\(h(Q_T)=\langle A,B,\gamma\rangle\).  Boolean operations on quantifiers
are pointwise, so \(T\) for \(Q_T\vee Q_{T'}\) is the pointwise join, and
\(h\) carries join, meet and complement to the coordinatewise operations of
\(\mathfrak C\).  Hence \(\mathcal A\) is closed under them (the
eventually constant sequences are closed under pointwise operations), is a
Boolean algebra, and \(h\) is an isomorphism.

For \(S_I\): \(T_1(n,\aleph_0)=\vee\) for all \(n\), \(T_1(\aleph_0,n)=
\wedge\), \(T_1(\aleph_0,\aleph_0)=\wedge\), giving \(X_1\).  For
\(S^{*}_I\): \(T_1^{*}(\mathfrak m,\mathfrak n)={\sim}T_1(\mathfrak n,
\mathfrak m)=\vee\) iff \(\mathfrak n\) is infinite, so \(T_1^{*}(n,\aleph_0)=\vee\),
\(T_1^{*}(\aleph_0,n)=\wedge\), \(T_1^{*}(\aleph_0,\aleph_0)=\vee\), giving
\(X_2\).  For \(S^0_I\) (with \(T_2\) as in
\cref{def:p03-w03-p03-042-misprint}): \(T_2(n,\aleph_0)=T_2(\aleph_0,n)=
\vee\) and \(T_2(\aleph_0,\aleph_0)=\wedge\), giving \(X_3\).

\emph{Item 3.}  \(\mathfrak I_0\) is closed under finite joins and contains
every element below one of its elements, so it is an ideal.  Its preimage
\(\mathcal A_0=h^{-1}[\mathfrak I_0]\) is an ideal because \(h\) is an
isomorphism.  For \(\langle A,B,\wedge\rangle\in\mathfrak I_0\) with \(A,B\)
finite, put
\[
 R=\bigvee_{m\in A}\overset{(m)}{\underset I\Sigma}\ \vee\
   \bigvee_{n\in B}\overset{(n)}{\underset I\Pi}
\]
(the empty join being \(\overset{(0)}{\underset I\Sigma}\wedge
{\sim}\overset{(0)}{\underset I\Sigma}\)).  At a pair \((k,\aleph_0)\),
\(\overset{(m)}{\Sigma}\) is \(\vee\) iff \(k=m\) and every
\(\overset{(n)}{\Pi}\) is \(\wedge\); at \((\aleph_0,k)\), every
\(\overset{(m)}{\Sigma}\) is \(\wedge\) and \(\overset{(n)}{\Pi}\) is \(\vee\)
iff \(k=n\); at \((\aleph_0,\aleph_0)\) all are \(\wedge\).  So
\(h(R)=\langle A,B,\wedge\rangle\) and every member of \(\mathcal A_0\) is
numerical.  Next, \(\mathfrak A/\mathfrak I\) has exactly two classes,
those of \(0\) (finite sets) and \(N\) (cofinite sets), and \(\{\wedge,
\vee\}\) is unchanged since the third factor of \(\mathfrak I_0\) is
\(\{\wedge\}\).  So every element of \(\mathfrak C\) is congruent modulo
\(\mathfrak I_0\) to exactly one \(\langle a,b,\gamma\rangle\) with
\(a,b\in\{0,N\}\), \(\gamma\in\{\wedge,\vee\}\).  These eight elements are
\(X_0=\langle0,0,\wedge\rangle\), \(X_1=\langle N,0,\wedge\rangle\),
\(X_2=\langle N,0,\vee\rangle\), \(X_3=\langle N,N,\wedge\rangle\),
\({\sim}X_0=\langle N,N,\vee\rangle\), \({\sim}X_1=\langle0,N,\vee\rangle\),
\({\sim}X_2=\langle0,N,\wedge\rangle\), \({\sim}X_3=\langle0,0,\vee\rangle\).

\emph{Item 4.}  Let \(Q_T\in\mathcal A\) and let \(Y\) be the one of the
eight elements congruent to \(h(Q_T)\), so \(h(Q_T)\triangle Y\in
\mathfrak I_0\), where \(\triangle\) is symmetric difference in the
Boolean algebra.  Let \(R=h^{-1}(h(Q_T)\triangle Y)\in\mathcal A_0\), a
numerical quantifier.  Then \(h^{-1}(Y)=Q_T\triangle R\).  If \(Y=X_0\),
\(h^{-1}(Y)\) is the constant quantifier with value \(\wedge\), so
\(Q_T=R\) is numerical.  If \(Y={\sim}X_0\), \(h^{-1}(Y)\) is the constant
\(\vee\) quantifier, equal to \(\overset{(0)}\Sigma\vee{\sim}
\overset{(0)}\Sigma\), so \(Q_T={\sim}R\) is numerical.  Otherwise
\(h^{-1}(Y)\) is one of \(S_I,S^{*}_I,S^0_I\) or their complements by
item~2, and \(h^{-1}(Y)=Q_T\triangle R\) is a Boolean polynomial in \(Q_T\)
and the numerical \(R\); taking a complement if needed, one of
\(S_I,S^{*}_I,S^0_I\) is such a polynomial.  A Boolean polynomial in
limited quantifiers is definable from them in the sense of
\cref{def:p03-w03-p03-042-definability} by the corresponding stroke
combination.
\end{proof}
\end{annexitem}

\begin{annexitem}{annex:p03-w03-p03-042-val-invariance}{thm:p03-w03-p03-042-val-invariance}{Free variables, isomorphism invariance and substitution}
In addition to the two stated parts, the proofs below use:
\emph{(3) substitution of a formula for a unary functional variable}.  Let
\(\mathrm F\) be unary, \(\mathrm W\) a formula and \(\mathrm z\) a
variable.  For a formula \(\mathrm Z\) whose bound individual variables do
not occur free in \(\mathrm W\), let \(\mathrm Z[\mathrm W/\mathrm F]\)
replace each atom \(\mathrm F(\mathrm x_i)\) by \(\mathrm W(\mathrm x_i/
\mathrm z)\), with \(\mathrm x_i\) not bound in \(\mathrm W\).  Then
\(\mathit{val}_{MI}(\mathrm Z[\mathrm W/\mathrm F])=\mathit{val}_{M'I}
(\mathrm Z)\), where \(M'\) agrees with \(M\) except that
\([\mathrm F]_{M'}=\underset{a\in I}{E}[\mathit{val}_{M(\mathrm z,a)I}
(\mathrm W)=\vee]\).
\begin{proof}
All three parts are proved by induction on \(\mathrm Z\), simultaneously
for all valuations.

(1) Atomic formulas: the value uses only \([\mathrm F]_M\) and the
\([\mathrm x_{i_r}]_M\) occurring.  Stroke: immediate.  For
\((\mathbf Q^j\mathrm x_i)\mathrm Z_1\): if \(M,N\) agree on the free
variables of it, then \(M(i,x),N(i,x)\) agree on those of \(\mathrm Z_1\)
for every \(x\), so the two sets in clause~4 coincide and so do the values.

(2) Atomic: \([\mathrm F]_{M'}(\varphi(a_1),\ldots)=([\mathrm F]_M)_\varphi
(\varphi(a_1),\ldots)=[\mathrm F]_M(a_1,\ldots)\), and \(\varphi\) preserves
and reflects equality.  Stroke: immediate.  Quantifier: \(M'(i,\varphi(x))\)
is the transport of \(M(i,x)\), so by induction the set \(X'\) defined in
\(I'\) from \(M'\) equals \(\varphi[X]\) for the set \(X\) defined from
\(M\); by the invariance clause of unlimited quantifiers
(\cref{def:p03-w03-p03-042-limited-quantifier}),
\(Q^j_{I'}(\varphi[X])=Q^j_I(X)\).

(3) Atoms \(\mathrm F(\mathrm x_i)\): the left side is
\(\mathit{val}_{MI}(\mathrm W(\mathrm x_i/\mathrm z))\), which by part~1 and
the substitution of a variable for a variable (a special case of the
clauses, as \(\mathrm x_i\) is free for \(\mathrm z\)) equals
\(\mathit{val}_{M(\mathrm z,[\mathrm x_i]_M)I}(\mathrm W)\), that is,
\(\vee\) iff \([\mathrm x_i]_M\in[\mathrm F]_{M'}\).  Other atoms and the
stroke: immediate.  Quantifier \((\mathbf Q^j\mathrm x_k)\mathrm Z_1\):
\(\mathrm x_k\) is not free in \(\mathrm W\), so by part~1 the set
\([\mathrm F]_{M(k,x)'}\) computed from \(M(k,x)\) equals
\([\mathrm F]_{M'}\), and the induction hypothesis applied to
\(M(k,x)\) gives the same defining set on both sides.
\end{proof}
\end{annexitem}

\begin{annexitem}{annex:p03-w03-p03-042-ring-lemma}{thm:p03-w03-p03-042-ring-lemma}{Rings satisfying the axioms}
\begin{proof}
(1) The map \(n\mapsto n\cdot1\) is a ring homomorphism, and it preserves
order since \(1\) is positive and \(\prec\) is compatible with addition;
being strictly increasing it is injective.  Let \(x\) be positive.  The
elements \(k\cdot1\) with \(k\geq1\) and \(k\cdot1\prec x\) lie in
\(Z_x\) and are distinct, so there are
finitely many; \(1\preceq x\) since \(1\) is the least positive element.
If \(x=1\) we are done; otherwise let \(k\geq1\) be largest with
\(k\cdot1\prec x\).  Then \(x\preceq(k+1)\cdot1\), and no element lies
strictly between \(k\cdot1\) and \((k+1)\cdot1\) by (d), so
\(x=(k+1)\cdot1\).  For negative \(x\), \(\ominus x\) is positive.  So the
map is onto.

(2) For an \(\Lang_A\)-formula \(\sigma\) of \cref{def:language-arithmetic}
define \(\sigma^{\mathbb Z}\).  Let \(\mathrm{One}(\mathrm u)\) be
\(\mathrm P(\mathrm u)\wedge(\forall\mathrm v)[\mathrm P(\mathrm v)\supset
{\sim}\mathrm F_4(\mathrm v,\mathrm u)]\) and \(\mathrm N(\mathrm x)\) be
\((\exists\mathrm u)[\mathrm F_1(\mathrm u)\wedge(\mathrm u=\mathrm x\vee
\mathrm F_4(\mathrm u,\mathrm x))]\).  For a term \(t\) and a variable
\(\mathrm y\) not in \(t\) define \(\tau_t(\mathrm y)\): \(\tau_0(\mathrm
y)=\mathrm F_1(\mathrm y)\); \(\tau_{\mathrm x}(\mathrm y)=(\mathrm y=
\mathrm x)\); \(\tau_{St}(\mathrm y)=(\exists\mathrm w)(\exists\mathrm u)
[\tau_t(\mathrm w)\wedge\mathrm{One}(\mathrm u)\wedge\mathrm F_2(\mathrm y,
\mathrm w,\mathrm u)]\); \(\tau_{t+r}(\mathrm y)=(\exists\mathrm w)(\exists
\mathrm w')[\tau_t(\mathrm w)\wedge\tau_r(\mathrm w')\wedge\mathrm F_2(
\mathrm y,\mathrm w,\mathrm w')]\), and \(\tau_{t\times r}\) likewise with
\(\mathrm F_3\); fresh variables are used throughout.  Put
\((t=r)^{\mathbb Z}=(\exists\mathrm y)[\tau_t(\mathrm y)\wedge\tau_r(
\mathrm y)]\), commute with the connectives, and
\(((\exists\mathrm x)\varphi)^{\mathbb Z}=(\exists\mathrm x)[\mathrm N(
\mathrm x)\wedge\varphi^{\mathbb Z}]\),
\(((\forall\mathrm x)\varphi)^{\mathbb Z}=(\forall\mathrm x)[\mathrm N(
\mathrm x)\supset\varphi^{\mathbb Z}]\).  At \(M_0\), \(\mathrm{One}\)
defines \(1\) and \(\mathrm N\) the non-negative integers.  By induction on
terms, for an assignment \(s\) of non-negative integers,
\(\tau_t(\mathrm y)\) holds at \(s\) and \(b\) iff \(b\) is the value of
\(t\) in \(\mathcal N\) at \(s\) (the value is non-negative,
and the ring operations of the integers restricted to non-negative
arguments are those of \(\mathcal N\), with \(S\) adding
\(1\)); by induction on formulas, \(\mathcal N\entails
\varphi[s]\) iff \(\varphi^{\mathbb Z}\) holds at \(M_0\) with the
variables assigned by \(s\), the quantifier steps using that \(\mathrm N\)
picks out the non-negative integers.  The map \(\sigma\mapsto\sigma^{
\mathbb Z}\) is computable.  If the set in (2) were recursively enumerable,
then listing it and comparing each listed formula with \(\sigma^{\mathbb
Z}\) would enumerate the true arithmetic sentences.  The set of their
codes would then be computably enumerable, so by the last clause of
\cref{thm:representability} it would be defined in \(\mathcal N\) by a
\(\Sigma^0_1\) formula; that formula would be an arithmetical truth
definition for the true \(\Lang_A\)-sentences, contradicting
\cref{thm:tarski-undefinability}.
\end{proof}
\end{annexitem}

\begin{annexitem}{annex:p03-w03-p03-042-theorem-two}{thm:p03-w03-p03-042-theorem-two}{Mostowski's Theorem~2}
\begin{proof}
Let \(\mathrm B_0\) be Mostowski's formula \cite[p.~17]{Mostowski1957}
\begin{multline*}
(\forall\mathrm x)\Bigl[\mathrm P(\mathrm x)\supset(\exists\mathrm y)\Bigl(
 \mathrm F_4(\mathrm x,\mathrm y)\wedge\bigl[(\mathbf Q\mathrm z)(\mathrm
 P(\mathrm z)\wedge\mathrm F_4(\mathrm z,\mathrm x))\not\equiv(\mathbf
 Q\mathrm z)(\mathrm P(\mathrm z)\wedge\mathrm F_4(\mathrm z,\mathrm
 y))\bigr]\\
 \wedge(\forall\mathrm t)\bigl\{\mathrm F_4(\mathrm x,\mathrm t)\wedge
 \mathrm F_4(\mathrm t,\mathrm y)\supset\bigl[(\mathbf Q\mathrm z)(\mathrm
 P(\mathrm z)\wedge\mathrm F_4(\mathrm z,\mathrm x))\equiv(\mathbf
 Q\mathrm z)(\mathrm P(\mathrm z)\wedge\mathrm F_4(\mathrm z,\mathrm t))
 \bigr]\bigr\}\Bigr)\Bigr],
\end{multline*}
which says that for every positive \(x\) there is a least \(y\succ x\)
with \(Q_I(Z_y)\neq Q_I(Z_x)\); let \(\mathrm A\) be \(\mathrm A_0\wedge
\mathrm B_0\).

(e) \emph{If \(M\) satisfies \(\mathrm A\) in \(I\), then \(I\) is
isomorphic to the ring of integers} \cite[p.~18]{Mostowski1957}.  By
\cref{thm:p03-w03-p03-042-ring-lemma}(1) it suffices that every \(Z_x\)
is finite.  Suppose \(Z_x\) is infinite for some positive \(x\)
(Mostowski takes \(\overline{Z_x}\) least; the argument below does not
need this).  The set of negative elements is mapped onto the positive ones
by \(\ominus\), and \(I\) is their union with \(\{\theta\}\); since the
positive elements include \(Z_x\), they are infinite, and with the axiom
of choice \(\overline I=\mathfrak p+1+\mathfrak p=\mathfrak p\), where
\(\mathfrak p\) is the power of the positive elements.  Every \(I-Z_u\)
with \(u\) positive contains all negative elements, so
\(\overline{I-Z_u}=\overline I\).  By \(\mathrm B_0\) take \(y\succ x\)
with \(Q_I(Z_x)\neq Q_I(Z_y)\), i.e.\ \(T(\overline{Z_x},\overline I)\neq
T(\overline{Z_y},\overline I)\), and with \(Q_I(Z_t)=Q_I(Z_x)\) for
\(x\prec t\prec y\).  If \(y=x\oplus1\), then \(Z_y=Z_x\cup\{x\}\) has the
power of the infinite \(Z_x\), contradicting the inequality.  So
\(x\prec t\prec y\) for \(t=y\ominus1\), by (d).  Now \(Z_y=Z_t\cup\{t\}\),
again by (d), so \(\overline{Z_t}=\overline{Z_y}\) (\(Z_y\supseteq Z_x\) is
infinite), and \(Q_I(Z_x)=Q_I(Z_t)=T(\overline{Z_t},\overline I)=
T(\overline{Z_y},\overline I)=Q_I(Z_y)\): a contradiction.

(f) \emph{The standard valuation \(M_0\) of the integers \(I_0\)
satisfies \(\mathrm A\)} \cite[p.~19]{Mostowski1957}.  \(\mathrm A_0\)
holds.  For positive \(x\), \(Z_x\) has \(x-1\) elements and infinite
complement, so \(Q_{I_0}(Z_x)=T(x-1,\aleph_0)=:a\).  By property (E) the
set of \(n\) with \(T(n,\aleph_0)\neq a\) is infinite; let \(n\) be its
least element \(\geq x-1\); then \(n\geq x\).  Put \(y=n+1\succ x\): then
\(Q(Z_y)=T(n,\aleph_0)\neq a\), and for \(x<t<y\), \(x\le t-1<n\), so
\(Q(Z_t)=T(t-1,\aleph_0)=a\).  So \(\mathrm B_0\) holds.

\emph{Reduction.}  Let \(\mathfrak X\) be the set of formulas of (S)
whose quantifiers are \(\exists,\forall\) and whose free variables are
among \(\mathrm F_1,\ldots,\mathrm F_4\).  For \(\mathrm X\in\mathfrak X\):
if \(\mathrm A\supset\mathrm X\) is true, then by (f)
\(\mathit{val}_{M_0I_0}(\mathrm X)=\vee\).  Conversely, let
\(\mathit{val}_{M_0I_0}(\mathrm X)=\vee\), and let \(M\) satisfy
\(\mathrm A\) in any \(I\).  By (e) there is an isomorphism \(\varphi\) of
\(I_0\) with its ring structure onto \(I\) with the structure given by
\(M\).  The transport of \(M_0\) along \(\varphi\) agrees with \(M\) on
\(\mathrm F_1,\ldots,\mathrm F_4\) (a propositional function with support
\(\{1,\ldots,k\}\) is fixed by its values on the first \(k\) coordinates),
so by \cref{thm:p03-w03-p03-042-val-invariance}(1),(2)
\(\mathit{val}_{MI}(\mathrm X)=\mathit{val}_{M_0I_0}(\mathrm X)=\vee\).
Hence \(\mathrm A\supset\mathrm X\) is true.  So \(\mathrm X\) holds at
\(M_0\) iff \(\mathrm A\supset\mathrm X\) is true.  \(\mathfrak X\) is
decidable and \(\mathrm X\mapsto\mathrm A\supset\mathrm X\) computable; if
the true formulas of (S) were recursively enumerable, so would be the set
of \cref{thm:p03-w03-p03-042-ring-lemma}(2), which is not.
\end{proof}
\end{annexitem}

\begin{annexitem}{annex:p03-w03-p03-042-theorem-three}{thm:p03-w03-p03-042-theorem-three}{Mostowski's Theorem~3}
\begin{proof}
Take \(\mathrm A=\mathrm A_0\wedge\mathrm C\), with \(\mathrm C\) one of
\((\forall\mathrm x)\{\mathrm P(\mathrm x)\supset(\mathbf S\mathrm z)
[\mathrm P(\mathrm z)\wedge\mathrm F_4(\mathrm z,\mathrm x)]\}\) and
\((\forall\mathrm x)\{\mathrm P(\mathrm x)\supset(\mathbf S^0\mathrm z)
[\mathrm P(\mathrm z)\wedge\mathrm F_4(\mathrm z,\mathrm x)]\}\)
\cite[p.~19]{Mostowski1957}.  If \(M\) satisfies \(\mathrm A\) in \(I\),
every \(Z_x\) (\(x\) positive) is finite: with \(S\) this is the meaning of
\(\mathrm C\); with \(S^0\), \(Z_x\) is finite or \(I-Z_x\) is finite,
and \(I-Z_x\) contains the infinitely many negative elements.  By
\cref{thm:p03-w03-p03-042-ring-lemma}(1), \(I\) is isomorphic to the
integers, and \(\mathrm A\) holds at \(M_0\) because each \(Z_x\) is
finite there.  The reduction in the proof of Theorem~2 now gives: \(\mathrm
X\in\mathfrak X\) holds at \(M_0\) iff \(\mathrm A\supset\mathrm X\) is
true iff \(\mathrm A\supset\mathrm X\) is true in a given denumerable
\(I\).  For the last equivalence: if \(\mathrm A\supset\mathrm X\) is true
in \(I\), transport \(M_0\) to \(I\) by a bijection; the transported
valuation satisfies \(\mathrm A\) by
\cref{thm:p03-w03-p03-042-val-invariance}(2), hence \(\mathrm X\), hence
\(\mathrm X\) holds at \(M_0\).  So neither the true formulas nor the
formulas true in \(I\) are recursively enumerable, by
\cref{thm:p03-w03-p03-042-ring-lemma}(2).
\end{proof}
\end{annexitem}

\begin{annexitem}{annex:p03-w03-p03-042-theorem-four}{thm:p03-w03-p03-042-theorem-four}{Mostowski's Theorem~4}
\begin{proof}
\emph{Sufficiency.}  Let each \(Q^j_I\) be a Boolean polynomial
\(\Phi_j\) in finitely many \(\overset{(m)}{\underset I\Sigma}\),
\(\overset{(n)}{\underset I\Pi}\).  For a formula \(\mathrm W\) and
\(\mathrm x\), let \(\mathrm E_m(\mathrm x,\mathrm W)\) be the formula with
\(\exists,\forall,=\) saying that exactly \(m\) values of \(\mathrm x\)
satisfy \(\mathrm W\) (for \(m=0\): \({\sim}(\exists\mathrm x)\mathrm W\);
for \(m\geq1\): there are pairwise distinct \(\mathrm y_1,\ldots,\mathrm
y_m\) satisfying \(\mathrm W\) and every \(\mathrm x\) satisfying \(\mathrm
W\) is one of them, with fresh \(\mathrm y_i\)).  Then
\(\overset{(m)}\Sigma\) corresponds to \(\mathrm E_m(\mathrm x,\mathrm W)\)
and \(\overset{(n)}\Pi\) to \(\mathrm E_n(\mathrm x,{\sim}\mathrm W)\).
Define \(\mathrm Z'\) by replacing, from the inside out, each
\((\mathbf Q^j\mathrm x)\mathrm W\) by \(\Phi_j\) applied to these
formulas of \(\mathrm W'\).  By induction and clause~4 of
\cref{def:p03-w03-p03-042-calculus-s}, \(\mathit{val}_{MI}(\mathrm Z)=
\mathit{val}_{MI}(\mathrm Z')\) for every \(I\)-valuation \(M\); the map
\(\mathrm Z\mapsto\mathrm Z'\) is computable from the finitely many
polynomials \(\Phi_j\).  Read the functional variables of \(\mathrm Z'\)
as relation symbols.  \(\mathrm Z'\) is true in \(I\) iff its universal
closure holds in every structure with universe \(I\), and by
\cref{thm:p03-w03-p03-042-val-invariance}(2) iff it holds in every
denumerable structure.  Let \(\Lambda\) be the recursive set of sentences
``there are at least \(k\) elements'', \(k\geq1\).  Every denumerable
structure is a model of \(\Lambda\); every model of \(\Lambda\) is
infinite and, by \cref{thm:downward-ls} for the finite signature of
\(\mathrm Z'\), a model of \(\Lambda\) in which the closure fails has a
countable elementary substructure in which it fails, which is infinite,
hence denumerable.  So \(\mathrm Z\) is true in \(I\) iff the closure of
\(\mathrm Z'\) is a consequence of \(\Lambda\), iff it is derivable from
\(\Lambda\) by \cref{cor:fol-completeness-countable}.  Derivations
from a recursive set can be listed, so the formulas true in \(I\) are
recursively enumerable.

\emph{Necessity.}  Suppose some \(Q^j_I\) is not numerical.

\emph{Case 1: a quantifier \(R_I\) definable in \(I\) in terms of
\(Q^1_I,\ldots,Q^s_I\) has property (E).}  Let \(\mathrm D\) define it,
with free variable \(\mathrm F\).  Form \(\mathrm B_0'\) from
\(\mathrm B_0\) by replacing each \((\mathbf Q\mathrm z)\mathrm W\) by
\(\mathrm D[\mathrm W/\mathrm F]\) (bound variables of \(\mathrm D\)
renamed).  By \cref{thm:p03-w03-p03-042-val-invariance}(3) and the defining
equivalence, \(\mathrm B_0'\) has in \(I\) the value that \(\mathrm B_0\)
has with \(\mathbf Q\) read as \(R_I\).  The arguments (e) and (f) of the
proof of Theorem~2 use only the quantifier on the one set \(I\), which is
denumerable; so every valuation on \(I\) satisfying \(\mathrm A_0\wedge
\mathrm B_0'\) yields a copy of the integers, and the transported standard
valuation satisfies it.  As in the proof of Theorem~3, \(\mathrm X\in
\mathfrak X\) holds at \(M_0\) iff \(\mathrm A_0\wedge\mathrm B_0'\supset
\mathrm X\) is true in \(I\), and the formulas true in \(I\) are not
recursively enumerable.

\emph{Case 2: no such quantifier exists.}  Let \(T\) determine \(Q^j_I\).
If \(T(n,\aleph_0)\) were not constant from some \(n\) on, both value sets
would be infinite and \(Q^j_I\), definable by \((\mathbf Q^j\mathrm x)
\mathrm F(\mathrm x)\), would have (E).  If \(T(\aleph_0,n)\) were not
eventually constant, then \(T^{*}(n,\aleph_0)={\sim}T(\aleph_0,n)\) would
not be, and the dual, definable by \({\sim}(\mathbf Q^j\mathrm x){\sim}
\mathrm F(\mathrm x)\) (Theorem~1), would have (E).  So
\(Q^j_I\in\mathcal A\) and is not numerical, and by
\cref{thm:p03-w03-p03-042-denumerable-classification}(4) one of
\(S_I,S^{*}_I,S^0_I\) is a Boolean polynomial in \(Q^j_I\) and numerical
quantifiers.  Numerical quantifiers are definable in \(I\) from
\(\exists,\forall\) by the formulas \(\mathrm E_m\) above, and \(S_I\) is
the dual of \(S^{*}_I\).  So \(S_I\) or \(S^0_I\) has a defining formula
\(\mathrm D\) in terms of \(Q^1_I,\ldots,Q^s_I\).  Substituting
\(\mathrm D\) in the formula \(\mathrm C\) of the proof of Theorem~3, as in
Case~1, gives a formula \(\mathrm A'\) of (S) such that \(\mathrm X\in
\mathfrak X\) holds at \(M_0\) iff \(\mathrm A'\supset\mathrm X\) is true
in \(I\); so the formulas true in \(I\) are not recursively enumerable.
\end{proof}
\end{annexitem}

\begin{annexitem}{annex:p03-w03-p03-042-theorem-five}{thm:p03-w03-p03-042-theorem-five}{Mostowski's Theorem~5}
\begin{proof}
Let \(\mathrm D\) define \(Q\) in terms of \(\exists,\forall,P\) and let
\(I\) be denumerable.  On \(I\), \(P_I\) is the constant quantifier with
value \(\vee\), so replacing every subformula \((\mathbf P\mathrm x)\mathrm
W\) of \(\mathrm D\) by \((\forall\mathrm x)(\mathrm x=\mathrm x)\) gives
\(\mathrm D_0\) with the same value at every \(I\)-valuation; \(\mathrm
D_0\) uses only \(\exists,\forall\).  Suppose \(Q_I=S_I\) (the case
\(S^0_I\) is the same with \(S^0\) in place of \(S\)).  Consider (S) with
the quantifiers \(\exists,\forall,S\).  For a formula \(\mathrm Y\), let
\(\mathrm Y^\dagger\) replace, from the inside out, each \((\mathbf
S\mathrm x)\mathrm W\) by \(\mathrm D_0[\mathrm W/\mathrm F]\).  By
\cref{thm:p03-w03-p03-042-val-invariance}(3), \(\mathrm Y\) and
\(\mathrm Y^\dagger\) have the same value at every \(I\)-valuation, and
\(\mathrm Y\mapsto\mathrm Y^\dagger\) is computable.  \(\mathrm Y^\dagger\)
uses only \(\exists,\forall\), which are numerical on \(I\)
(\(\exists_I={\sim}\overset{(0)}{\underset I\Sigma}\),
\(\forall_I=\overset{(0)}{\underset I\Pi}\)), so by the sufficiency half
of Theorem~4 the formulas \(\mathrm Y^\dagger\) true in \(I\) form a
recursively enumerable set.  Then so would the formulas \(\mathrm Y\) with
\(\exists,\forall,S\) true in \(I\), contradicting Theorem~3.
\end{proof}
\end{annexitem}

\begin{annexitem}{annex:p03-w03-p03-042-transfer}{thm:p03-w03-p03-042-transfer}{Determining functions across sets}
\begin{proof}
(1) Let \(\varphi\) map \(I\) one-one onto \(I'\), and let
\(\mathfrak m+\mathfrak n=\overline I\).  By
\cref{thm:p03-w03-p03-042-theorem-one} there is a one-argument \(F\) on
\(I\) with \(\overline{F^{-1}(\vee)}=\mathfrak m\) and
\(\overline{F^{-1}(\wedge)}=\mathfrak n\).  Then \(F_\varphi^{-1}(\vee)=
\varphi[F^{-1}(\vee)]\) and \(F_\varphi^{-1}(\wedge)=\varphi[F^{-1}(\wedge)]\)
have the same powers, so \(T_I(\mathfrak m,\mathfrak n)=Q_I(F)=
Q_{I'}(F_\varphi)=T_{I'}(\mathfrak m,\mathfrak n)\) by the invariance
clause of \cref{def:p03-w03-p03-042-limited-quantifier}.  The pairs with
sum \(\overline I\) and with sum \(\overline{I'}\) are the same pairs.

(2) Boolean operations are pointwise, so \(\Phi(R^1_I,\ldots,R^r_I)(F)=
\Phi(R^1_I(F),\ldots,R^r_I(F))=\Phi(T^1_I(\mathfrak m,\mathfrak n),\ldots,
T^r_I(\mathfrak m,\mathfrak n))\) with \(\mathfrak m,\mathfrak n\) the
powers of \(F^{-1}(\vee),F^{-1}(\wedge)\); by the uniqueness clause of
\cref{thm:p03-w03-p03-042-theorem-one} this pointwise \(\Phi\) is the
determining function.  Each equation of
\cref{def:p03-w03-p03-042-infinite-powers} holds for every \(T^i\), so it
holds after applying \(\Phi\) to both sides.

(3) Let \(T,T'\) determine \(R,R'\).  \(R_I=R'_I\) gives \(T_I=T'_I\) by
uniqueness in \cref{thm:p03-w03-p03-042-theorem-one}.  Let \(J\) be
infinite and \(\mathfrak m+\mathfrak n=\overline J\).  If \(\mathfrak m=k\)
is finite, then \(\mathfrak n\) is infinite and \(\mathfrak n=k+\mathfrak
n=\overline J\), so the first clause of
\cref{def:p03-w03-p03-042-infinite-powers} for \(R\) and \(R'\) gives
\(T_J(k,\overline J)=T_I(k,\aleph_0)=T'_I(k,\aleph_0)=T'_J(k,\overline J)\).
If \(\mathfrak n\) is finite, the second clause gives the same.  If both
are infinite, the third clause, applied to \(J\) and to \(I\) with the pair
\((\aleph_0,\aleph_0)\), gives \(T_J(\mathfrak m,\mathfrak n)=
T_I(\aleph_0,\aleph_0)=T'_I(\aleph_0,\aleph_0)=T'_J(\mathfrak m,\mathfrak
n)\).  So \(T_J=T'_J\) and \(R_J=R'_J\).
\end{proof}
\end{annexitem}

\begin{annexitem}{annex:p03-w03-p03-042-lemma-a}{thm:p03-w03-p03-042-lemma-a}{Mostowski's Lemma~(a) of \S5}
\begin{proof}
The proof of \cref{thm:p03-w03-p03-042-val-invariance}(1) applies to
improper formulas without change, since a constant has the same value
\([\sigma_x]_M=x\) at every valuation.  Write \(c=[\mathrm x_j]_M\) and
\(\mathrm A^c=\mathrm A(\sigma_c/\mathrm x_j)\) (free occurrences only).
By induction on \(\mathrm A\), for all valuations \(M\):
\(\mathit{val}_{MI}(\mathrm A)=\mathit{val}_{MI}(\mathrm A^{[\mathrm
x_j]_M})\).  Atomic formulas: an occurrence of \(\mathrm x_j\) and of
\(\sigma_c\) both take the value \(c\).  Stroke: immediate.  Quantifier
\((\mathbf Q^h\mathrm x_i)\mathrm A_1\): if \(i=j\), \(\mathrm x_j\) is not
free and \(\mathrm A^c=\mathrm A\).  If \(i\neq j\), then
\([\mathrm x_j]_{M(i,x)}=c\) for every \(x\in I\), so by the induction
hypothesis at \(M(i,x)\) the sets \(\underset{x\in I}{E}[\mathit{val}_{M(i,x)I}
(\mathrm A_1)=\vee]\) and \(\underset{x\in I}{E}[\mathit{val}_{M(i,x)I}
(\mathrm A_1^c)=\vee]\) coincide, and clause~4 of
\cref{def:p03-w03-p03-042-calculus-s} gives equal values.

For the first statement apply this successively to \(\mathrm x_1,\ldots,
\mathrm x_l\) (the value of \(\mathrm x_{j'}\) at \(M\) is unchanged by the
earlier substitutions): \(\mathit{val}_{MI}(\mathrm A)=\mathit{val}_{MI}
(\mathrm A')\).  \(\mathrm A'\) has no free individual variables, so by
part~(1) its value depends only on \([\mathrm F_1]_M,\ldots,[\mathrm
F_k]_M\); hence \(\mathit{val}_{MI}(\mathrm A)=\vee\) iff these functions
satisfy \(\mathrm A'\) in \(I\).
\end{proof}
\end{annexitem}

\begin{annexitem}{annex:p03-w03-p03-042-lemma-b}{thm:p03-w03-p03-042-lemma-b}{Mostowski's Lemma~(b) of \S5}
\begin{proof}
Following Mostowski \cite[pp.~22--23]{Mostowski1957}.  Since every
\([\mathrm x_i]_M\) lies in \(I_0\), \(M_0\) gives each \(\mathrm x_i\) the
same value as \(M\), and for \(x_1,x_2,\ldots\in I_0\)
\[
 [\mathrm F_j]_{M_0}(x_1,x_2,\ldots)=[\mathrm F_j]_M(x_1,x_2,\ldots).
 \tag{3}
\]
Induction on the proper formula \(\mathrm V\), for all admissible
substitutions at once.

\emph{Atomic.}  \(\mathrm V'\) is \(\mathrm F_j(t_1,\ldots,t_k)\) or
\(t_1=t_2\), each \(t_r\) a variable or a constant \(\sigma_a\) with
\(a\in I_0\).  At \(M\) and at \(M_0\) every \(t_r\) has the same value,
which lies in \(I_0\); by (3) the two values of \(\mathrm V'\) agree.

\emph{Stroke.}  \(\mathrm V'=\mathrm V_1'|\mathrm V_2'\) with
\(\mathrm V_1',\mathrm V_2'\) obtained by the same substitution; apply the
induction hypothesis to both.

\emph{Quantifier.}  \(\mathrm V=(\mathbf Q^h\mathrm x_i)\mathrm V_1\) and
\(\mathrm V'=(\mathbf Q^h\mathrm x_i)\mathrm V_1'\).  Let \(\mathrm y_1,
\ldots,\mathrm y_k\) be the free individual variables of \(\mathrm V'\) and
substitute \(\sigma_{[\mathrm y_r]_M}\) for \(\mathrm y_r\); the result is
a formula \((\mathbf Q^h\mathrm x_i)\mathrm W\) without free individual
variables, obtained from \((\mathbf Q^h\mathrm x_i)\mathrm V_1\in
\mathfrak U\) by substituting constants with indices in \(I_0\) for all its
free individual variables.  There are finitely many such constants; since
\(I_1\subseteq I_2\subseteq\cdots\), all of them have indices in one
\(I_k\), and \((\mathbf Q^h\mathrm x_i)\mathrm W\) is a term \(\mathrm
V_j\) of the list formed at stage \(k\) of \cref{def:p03-w03-p03-042-hull}.
By \cref{thm:p03-w03-p03-042-lemma-a}, and since \([\mathrm
y_r]_{M_0}=[\mathrm y_r]_M\),
\[
 \mathit{val}_{MI}(\mathrm V')=\mathit{val}_{MI}((\mathbf Q^h\mathrm x_i)
 \mathrm W),\qquad
 \mathit{val}_{M_0I_0}(\mathrm V')=\mathit{val}_{M_0I_0}((\mathbf Q^h
 \mathrm x_i)\mathrm W),
\]
so it suffices to prove
\[
 \mathit{val}_{MI}((\mathbf Q^h\mathrm x_i)\mathrm W)=
 \mathit{val}_{M_0I_0}((\mathbf Q^h\mathrm x_i)\mathrm W).\tag{6}
\]
By clause~4, the left side is \(Q^h_I\) of the set of \(x\in I\) with
\(\mathit{val}_{M(i,x)I}(\mathrm W)=\vee\); by
\cref{thm:p03-w03-p03-042-lemma-a} this value equals
\(\mathit{val}_{M(i,x)I}(\mathrm W(\sigma_x))\), and since
\(\mathrm W(\sigma_x)\) has no free individual variables it equals
\(\mathit{val}_{MI}(\mathrm W(\sigma_x))\), which depends only on the
\([\mathrm F_j]_M=[\mathrm F_j]_{\overline M}\); so it equals
\(\mathit{val}_{\overline MI}(\mathrm W(\sigma_x))\), for every \(x\in I\).
With \(T^h_I\) determining \(Q^h_I\) and \(J_1,J_2\) the sets formed for
\(\mathrm V_j\) at stage \(k\),
\[
 \mathit{val}_{MI}((\mathbf Q^h\mathrm x_i)\mathrm W)=T^h_I(\mathfrak m_1,
 \mathfrak m_2),\qquad\mathfrak m_1=\overline{J_1},\ \mathfrak m_2=
 \overline{J_2}.\tag{7}
\]
In the same way
\[
 \mathit{val}_{M_0I_0}((\mathbf Q^h\mathrm x_i)\mathrm W)=T^h_{I_0}
 (\mathfrak m^0_1,\mathfrak m^0_2),\tag{8}
\]
with \(\mathfrak m^0_1,\mathfrak m^0_2\) the powers of the sets of
\(a\in I_0\) at which \(\mathit{val}_{M_0I_0}(\mathrm W(\sigma_a))\) is
\(\vee\), resp.\ \(\wedge\).  For \(a\in I_0\), \(\mathrm W(\sigma_a)\)
results from \(\mathrm V_1\) by substituting constants with indices in
\(I_0\) for all free individual variables of \(\mathrm V_1\), so the
induction hypothesis gives \(\mathit{val}_{M_0I_0}(\mathrm W(\sigma_a))=
\mathit{val}_{MI}(\mathrm W(\sigma_a))=\mathit{val}_{\overline MI}(\mathrm
W(\sigma_a))\).  Hence \(\mathfrak m^0_1=\overline{J_1\cap I_0}\) and
\(\mathfrak m^0_2=\overline{J_2\cap I_0}\).

Now \(I_{kj}\subseteq I_{k+1}\subseteq I_0\) and \(I_0\) is denumerable.
If \(\mathfrak m_1,\mathfrak m_2\) are both infinite, \(J_1\cap I_0\) and
\(J_2\cap I_0\) are infinite, so \(\mathfrak m^0_1=\mathfrak m^0_2=
\aleph_0\), and the third clause of
\cref{def:p03-w03-p03-042-infinite-powers} gives \(T^h_I(\mathfrak m_1,
\mathfrak m_2)=T^h_{I_0}(\aleph_0,\aleph_0)\).  If \(\mathfrak m_1=n\) is
finite, then \(J_1\subseteq I_0\), so \(\mathfrak m^0_1=n\), and
\(J_2\cap I_0\) is infinite, so \(\mathfrak m^0_2=\aleph_0=\overline{I_0}\);
moreover \(\mathfrak m_2=\overline I\) because \(I\) is infinite; the
first clause gives \(T^h_I(n,\overline I)=T^h_{I_0}(n,\overline{I_0})\).
The case of finite \(\mathfrak m_2\) is symmetric, with the second clause.
Both cannot be finite since \(J_1\cup J_2=I\) is infinite.  In all cases
the right sides of (7) and (8) agree, which proves (6).
\end{proof}
\end{annexitem}

\begin{annexitem}{annex:p03-w03-p03-042-theorem-six}{thm:p03-w03-p03-042-theorem-six}{Mostowski's Theorem~6}
\begin{proof}
Let \(\mathrm Z\) be closed and \(\mathit{val}_{\overline MI}(\mathrm
Z)=\vee\) for an infinite \(I\) and an \(I\)-valuation \(\overline M\).
Build \(I_0\) from \(I\) and \(\overline M\)
(\cref{def:p03-w03-p03-042-hull}).  Let \(M\) agree with \(\overline M\)
on all functional variables and give every individual variable a value in
\(I_0\).  \(\mathrm Z\) has no free individual variables, so
\(\mathit{val}_{MI}(\mathrm Z)=\mathit{val}_{\overline MI}(\mathrm Z)=\vee\)
by \cref{thm:p03-w03-p03-042-val-invariance}(1).  \(M\) satisfies the
hypotheses of \cref{thm:p03-w03-p03-042-lemma-b}; applied to
\(\mathrm V=\mathrm Z\) with the empty substitution it gives
\(\mathit{val}_{M_0I_0}(\mathrm Z)=\vee\).  So \(\mathrm Z\) is
satisfiable in the denumerable set \(I_0\) \cite[p.~23]{Mostowski1957}.
\end{proof}
\end{annexitem}

\begin{annexitem}{annex:p03-w03-p03-042-theorem-seven}{thm:p03-w03-p03-042-theorem-seven}{Mostowski's Theorem~7}
\begin{proof}
Let \(Q=Q^h\) distinguish infinite powers.  By
\cref{thm:p03-w03-p03-042-transfer}(1) write \(T_{\mathfrak p}\) for the
function determining \(Q_I\) when \(\overline I=\mathfrak p\).  For a
formula \(\mathrm H\), \(\neg_\varepsilon(\mathbf Q\mathrm u)\mathrm H\)
denotes \((\mathbf Q\mathrm u)\mathrm H\) if \(\varepsilon=\wedge\) and
\({\sim}(\mathbf Q\mathrm u)\mathrm H\) if \(\varepsilon=\vee\); its value
is \(\vee\) iff \(Q\) gives the set defined by \(\mathrm H\) the value
\({\sim}\varepsilon\).  Abbreviations (\(\exists,\forall\) are among the
quantifiers):
\begin{gather*}
\mathrm{Ord}:\ {\sim}(\exists\mathrm x)\mathrm G(\mathrm x,\mathrm x)\wedge
(\forall\mathrm x,\mathrm y)[{\sim}\mathrm G(\mathrm x,\mathrm y)\vee
{\sim}\mathrm G(\mathrm y,\mathrm x)]\wedge(\forall\mathrm x,\mathrm y,
\mathrm z)[\mathrm G(\mathrm x,\mathrm y)\wedge\mathrm G(\mathrm y,\mathrm
z)\supset\mathrm G(\mathrm x,\mathrm z)]\\
\wedge(\forall\mathrm x,\mathrm y)[\mathrm G(\mathrm x,\mathrm y)\vee
\mathrm x=\mathrm y\vee\mathrm G(\mathrm y,\mathrm x)],\\
\mathrm{Lim}(\mathrm y):\ (\forall\mathrm s)\{\mathrm G(\mathrm s,\mathrm
y)\supset(\exists\mathrm t)[\mathrm G(\mathrm s,\mathrm t)\wedge\mathrm
G(\mathrm t,\mathrm y)]\}.
\end{gather*}
\(\mathrm{Ord}\) says that \(\mathrm G\) is a strict linear order of the
universe, written \(<\).  Two facts:

(L) \emph{If \(\mathrm{Ord}\) holds, \(r<y\) and \(\mathrm{Lim}(y)\), then
infinitely many elements lie strictly between \(r\) and \(y\).}  Put
\(t_0=r\) and, given \(t_k<y\), choose \(t_{k+1}\) with \(t_k<t_{k+1}<y\)
(dependent choice); the \(t_k\) increase strictly, so they are distinct.

(W) \emph{A set of infinite power \(\mathfrak p\) has a well ordering of
order type at least \(\omega+1\); in such an ordering the element in
position \(\omega\) has a predecessor and satisfies \(\mathrm{Lim}\).}  By
the axiom of choice the set is equipotent with an infinite ordinal
\(\kappa\), hence with \(\kappa+1\); transport the ordering.  The element
\(e_\omega\) in position \(\omega\) is preceded by \(e_0\), and any
\(s<e_\omega\) is some \(e_k\), \(k<\omega\), with \(e_k<e_{k+1}<e_\omega\).

\emph{Case (a): for some \(n\), \(f(\mathfrak p)=T_{\mathfrak p}(n,
\mathfrak p)\) is not constant on infinite cardinals.}  Mostowski takes
\(\mathfrak m_2\) least with a counterexample below it
\cite[p.~24]{Mostowski1957}; the argument needs only an infinite
\(\mathfrak m_2\) with \(f(\mathfrak m_2)\neq f(\aleph_0)\), which exists
because \(f\) is not constant, and then \(\mathfrak m_2>\aleph_0\).  Let
\(\varepsilon=f(\aleph_0)\) and let \(\mathrm A_{(a)}\) be
\begin{multline*}
\mathrm{Ord}\wedge(\exists\mathrm x_0)\Bigl\{(\exists\mathrm x_1,\ldots,
\mathrm x_n)\Bigl(\mathbf K_{i<j}(\mathrm x_i\neq\mathrm x_j)\wedge
(\forall\mathrm z)[\mathrm G(\mathrm z,\mathrm x_0)\equiv(\mathrm z=
\mathrm x_1\vee\cdots\vee\mathrm z=\mathrm x_n)]\Bigr)\\
\wedge(\exists\mathrm y)\bigl(\mathrm G(\mathrm x_0,\mathrm y)\wedge
\mathrm{Lim}(\mathrm y)\bigr)\wedge\neg_\varepsilon(\mathbf Q\mathrm u)
\mathrm G(\mathrm u,\mathrm x_0)\Bigr\},
\end{multline*}
where \(\mathbf K_{i<j}(\mathrm x_i\neq\mathrm x_j)\) is the conjunction of
all \(\mathrm x_i\neq\mathrm x_j\), \(i<j\le n\)
\cite[p.~24]{Mostowski1957}, and for \(n=0\) the first conjunct is
\((\forall\mathrm z){\sim}\mathrm G(\mathrm z,\mathrm x_0)\).  Mostowski
assumes \(\varepsilon=\wedge\) without loss of generality; the operator
\(\neg_\varepsilon\) makes that step explicit.

\emph{Satisfiable in a set \(I_2\) of power \(\mathfrak m_2\).}  Well order
\(I_2\) (axiom of choice).  As \(I_2\) is not denumerable, the order type
exceeds \(\omega\), so positions \(0,\ldots,n\) and \(\omega\) are
occupied.  Let \([\mathrm G]_M(y_1,y_2,\ldots)=\vee\) iff \(y_1\) precedes
\(y_2\) \cite[p.~24]{Mostowski1957}.  Take \(\mathrm x_0\) in position
\(n\), \(\mathrm x_1,\ldots,\mathrm x_n\) in positions \(0,\ldots,n-1\),
and \(\mathrm y\) in position \(\omega\), which satisfies \(\mathrm{Lim}\)
as in (W).  The set \(S\) of predecessors of \(\mathrm x_0\) has \(n\)
elements and \(I_2-S\) has power \(\mathfrak m_2\), so
\(Q_{I_2}(S)=f(\mathfrak m_2)={\sim}\varepsilon\), and the last conjunct
holds.

\emph{Not satisfiable in a denumerable \(I\).}  Otherwise \(<\) orders
\(I\) and some \(x_0\) has exactly \(n\) predecessors; their set \(S\) is
finite and \(I-S\) has power \(\aleph_0\), so \(Q_I(S)=T_{\aleph_0}(n,
\aleph_0)=\varepsilon\), and the last conjunct fails
\cite[pp.~24--25]{Mostowski1957}.

\emph{Case (b): for some \(n\), \(g(\mathfrak p)=T_{\mathfrak p}(\mathfrak
p,n)\) is not constant.}  Mostowski says this case is similar
\cite[p.~25]{Mostowski1957}.  Take \(\mathfrak m_2\) with \(g(\mathfrak
m_2)\neq g(\aleph_0)=\varepsilon\), and form \(\mathrm A_{(b)}\) from
\(\mathrm A_{(a)}\) by replacing the last conjunct with
\(\neg_\varepsilon(\mathbf Q\mathrm u){\sim}\mathrm G(\mathrm u,\mathrm
x_0)\).  The set defined by \({\sim}\mathrm G(\mathrm u,\mathrm x_0)\) is
the complement of the set \(S\) of predecessors of \(\mathrm x_0\); \(Q\)
gives it the value \(T_{\overline I}(\overline{I-S},\overline S)\).  In the
well ordering of \(I_2\) above this is \(g(\mathfrak m_2)={\sim}
\varepsilon\); in a denumerable \(I\) with \(\overline S=n\) it is
\(g(\aleph_0)=\varepsilon\).  The two halves of the argument are otherwise
those of case~(a).

\emph{Case (c): neither (a) nor (b) applies.}  Then the third clause of
\cref{def:p03-w03-p03-042-infinite-powers} fails: there are infinite
\(\mathfrak m_i,\mathfrak n_i\) with \(\mathfrak m_i+\mathfrak n_i=
\mathfrak p_i\) and \(T_{\mathfrak p_1}(\mathfrak m_1,\mathfrak n_1)\neq
T_{\mathfrak p_2}(\mathfrak m_2,\mathfrak n_2)\).  Let \(\varepsilon=
T_{\aleph_0}(\aleph_0,\aleph_0)\).  One of the two differing values is
not \(\varepsilon\); rename so that \(T_{\mathfrak p_2}(\mathfrak m_2,
\mathfrak n_2)={\sim}\varepsilon\).  Then \(\mathfrak p_2>\aleph_0\), since
\((\aleph_0,\aleph_0)\) is the only pair of infinite cardinals with sum
\(\aleph_0\).  Mostowski chooses \(\mathfrak p_2\) as small as possible and
derives \(T_{\aleph_0}(\aleph_0,\aleph_0)\neq T_{\mathfrak p_2}(\mathfrak
m_2,\mathfrak n_2)\) from minimality \cite[p.~25]{Mostowski1957}; the
renaming above gives this directly.  Let \(\mathrm A_{(c)}\) be
\begin{multline*}
\mathrm{Ord}\wedge(\exists\mathrm x_0)\Bigl\{\neg_\varepsilon(\mathbf
Q\mathrm u)\mathrm G(\mathrm u,\mathrm x_0)\wedge(\exists\mathrm y_1)
\bigl[\mathrm G(\mathrm y_1,\mathrm x_0)\wedge(\exists\mathrm r)\mathrm
G(\mathrm r,\mathrm y_1)\wedge\mathrm{Lim}(\mathrm y_1)\bigr]\\
\wedge(\exists\mathrm y_2)\bigl[\mathrm G(\mathrm x_0,\mathrm y_2)\wedge
\mathrm{Lim}(\mathrm y_2)\bigr]\Bigr\},
\end{multline*}
which is Mostowski's formula with the conjunct \((\exists\mathrm r)\mathrm
G(\mathrm r,\mathrm y_1)\) of \cref{def:p03-w03-p03-042-limit-point}.

\emph{Satisfiable in a set \(I_2\) of power \(\mathfrak p_2\).}  Split
\(I_2\) into disjoint \(S,R\) of powers \(\mathfrak m_2,\mathfrak n_2\).
By (W) well order \(S\) and \(R\) each with order type at least
\(\omega+1\), and order \(I_2\) by putting \(S\) before \(R\); let
\(\mathrm G\) express this order.  Take \(\mathrm x_0\) least in \(R\), so
that its predecessors form \(S\); \(\mathrm y_1\) in position \(\omega\) of
\(S\), with \(\mathrm r\) least in \(S\); \(\mathrm y_2\) in position
\(\omega\) of \(R\).  \(\mathrm{Lim}(\mathrm y_1)\) holds by (W).  For
\(\mathrm{Lim}(\mathrm y_2)\): if \(s<\mathrm y_2\) and \(s\in S\), take
\(t=\mathrm x_0\); if \(s\in R\), \(s\) is in a finite position \(k\) of
\(R\) and the element in position \(k+1\) lies between.  Finally
\(Q_{I_2}(S)=T_{\mathfrak p_2}(\mathfrak m_2,\mathfrak n_2)=
{\sim}\varepsilon\).

\emph{Not satisfiable in a denumerable \(I\).}  Otherwise let \(S\) be the
set of predecessors of \(x_0\).  It contains \(y_1\) and, by (L) applied to
\(r<y_1\), infinitely many elements below \(y_1\).  Its complement contains
\(x_0\) and, by (L) applied to \(x_0<y_2\), infinitely many elements.  As
\(I\) is denumerable both have power \(\aleph_0\), so
\(Q_I(S)=T_{\aleph_0}(\aleph_0,\aleph_0)=\varepsilon\), and the first
conjunct fails \cite[p.~25]{Mostowski1957}.

In each case the formula is closed and satisfiable in the non-denumerable
set \(I_2\) but in no denumerable set.
\end{proof}
\end{annexitem}

\begin{annexitem}{annex:p03-w03-p03-042-theorem-eight}{thm:p03-w03-p03-042-theorem-eight}{Mostowski's Theorem~8}
\begin{proof}
Following Mostowski \cite[pp.~26--27]{Mostowski1957}, with his
computation (11)--(14) replaced by \cref{thm:p03-w03-p03-042-transfer}.

\emph{Step 1.  None of \(Q^1,\ldots,Q^s\) distinguishes infinite powers.}
Otherwise \cref{thm:p03-w03-p03-042-theorem-seven} gives a closed formula
satisfiable in a non-denumerable, hence infinite, set but not in a
denumerable one, against the hypothesis.

\emph{Step 2.  If \(I\) is denumerable, no quantifier \(R_I\) definable in
\(I\) in terms of \(Q^1_I,\ldots,Q^s_I\) has property (E).}  Let
\(\mathrm D\) define such an \(R_I\) with (E), and let \(\mathrm A_0\wedge
\mathrm B_0'\) be formed as in Case~1 of the proof of
\cref{thm:p03-w03-p03-042-theorem-four}; as shown there, it is satisfiable
in \(I\).  Let \(J\) be any set and \(M\) a \(J\)-valuation satisfying it.
By \cref{thm:p03-w03-p03-042-val-invariance}(1),(2), the map sending a
one-argument \(X\) on \(J\) to the value of \(\mathrm D\) at a valuation
giving \(\mathrm F\) the value \(X\) is well defined and invariant under
permutations of \(J\), so it is a quantifier \(R'_J\) limited to \(J\); by
part~(3), \(\mathrm B_0'\) has at \(M\) the value of \(\mathrm B_0\) with
\(\mathbf Q\) read as \(R'_J\).  Argument (e) in the proof of
\cref{thm:p03-w03-p03-042-theorem-two} uses nothing about the quantifier
beyond its determining function on the one set, so \(J\) with the
structure given by \(M\) is isomorphic to the ring of integers and \(J\) is
denumerable.  So \(\mathrm A_0\wedge\mathrm B_0'\) is satisfiable in the
infinite set \(I\) but in no non-denumerable set, against the hypothesis.

\emph{Step 3.  Infinite \(J\).}  Fix \(j\) and a denumerable \(I\).
\(Q^j_I\) and its dual, definable by \({\sim}(\mathbf Q^j\mathrm x)
{\sim}\mathrm F(\mathrm x)\), lack (E) by Step~2, so, as in Case~2 of the
proof of \cref{thm:p03-w03-p03-042-theorem-four}, \(Q^j_I\) lies in the
algebra \(\mathcal A\) of
\cref{thm:p03-w03-p03-042-denumerable-classification}.

Suppose \(Q^j_I\) is not numerical.  By
\cref{thm:p03-w03-p03-042-denumerable-classification}(4) one of
\(U_I\in\{S_I,S^{*}_I,S^0_I\}\) equals \(\Phi(Q^j_I,N^1_I,\ldots,N^r_I)\)
for a Boolean polynomial \(\Phi\) and numerical quantifiers \(N^i\); this
is Mostowski's (9) \cite[p.~26]{Mostowski1957}.  Let \(R\) be the unlimited
quantifier \(\Phi(Q^j,N^1,\ldots,N^r)\).  \(Q^j\) (Step~1) and the
\(\overset{(n)}{\Sigma},\overset{(m)}{\Pi}\) do not distinguish infinite
powers, so neither do \(R\) nor \(U\)
(\cref{thm:p03-w03-p03-042-transfer}(2) and the remark after
\cref{def:p03-w03-p03-042-infinite-powers}).  Since \(R_I=U_I\),
\cref{thm:p03-w03-p03-042-transfer}(3) gives \(R_J=U_J\) for every infinite
\(J\): Mostowski's conclusion that (9) holds for every infinite set
\cite[p.~27]{Mostowski1957}.  Let \(\mathrm D_R\) be the stroke
combination of \((\mathbf Q^j\mathrm x)\mathrm F(\mathrm x)\) and the
formulas \(\mathrm E_m(\mathrm x,\pm\mathrm F(\mathrm x))\) of the proof of
\cref{thm:p03-w03-p03-042-theorem-four} that corresponds to \(\Phi\); it
defines \(R_J\) in every \(J\).  If \(U=S^{*}\), then \(S_J(F)={\sim}
S^{*}_J({\sim}F)\), and \({\sim}\mathrm D_R[{\sim}\mathrm F(\mathrm z)/
\mathrm F]\) defines \(S_J\) in every infinite \(J\).  So \(S\) or \(S^0\)
has a formula \(\mathrm D\) in \(\mathbf Q^j,\exists,\forall\) defining it
in every infinite set.  Substitute \(\mathrm D\) for \(\mathbf S\) (resp.\
\(\mathbf S^0\)) in the formula \(\mathrm A_0\wedge\mathrm C\) of the proof
of \cref{thm:p03-w03-p03-042-theorem-three}, as in Step~2.  \(\mathrm
A_0\) forces the universe to be infinite, so at every valuation
satisfying \(\mathrm A_0\) the result has the value of \(\mathrm A_0\wedge
\mathrm C\); by that proof it is satisfiable in the integers and every set
in which it is satisfiable is denumerable, against the hypothesis
\cite[p.~27]{Mostowski1957}.

Hence \(Q^j_I=\Psi(N^1_I,\ldots,N^q_I)\) for a Boolean polynomial \(\Psi\)
and quantifiers \(N^i\) among the \(\overset{(n)}\Sigma,\overset{(m)}\Pi\).
Both \(Q^j\) and \(\Psi(N^1,\ldots,N^q)\) do not distinguish infinite
powers, so by \cref{thm:p03-w03-p03-042-transfer}(3)
\(Q^j_J=\Psi(N^1_J,\ldots,N^q_J)\) for every infinite \(J\), with the same
\(\Psi\).  The stroke combination of the formulas \(\mathrm E_m\)
corresponding to \(\Psi\) uses only \(\exists,\forall\) and defines
\(Q^j_J\) in \(J\).

\emph{Step 4.  Finite \(J\)} (Mostowski: evident
\cite[p.~27]{Mostowski1957}).  Let \(J\) have \(k\) elements and \(T\)
determine \(Q^j_J\).  The pairs with sum \(k\) are \((m,k-m)\),
\(m=0,\ldots,k\).  With \(A=\{m\le k:T(m,k-m)=\vee\}\),
\(Q^j_J(F)=\vee\) iff \(\overline{F^{-1}(\vee)}\in A\), so \(Q^j_J\) is
defined in \(J\) by the disjunction of the \(\mathrm E_m(\mathrm x,\mathrm
F(\mathrm x))\), \(m\in A\), or by \({\sim}(\forall\mathrm x)(\mathrm x=
\mathrm x)\) if \(A\) is empty.
\end{proof}
\end{annexitem}

\begin{annexitem}{annex:p03-w03-p03-042-theorem-eight-boundary}{thm:p03-w03-p03-042-theorem-eight-boundary}{The boundary example}
\begin{proof}
\emph{\(Q\) is an unlimited quantifier.}  Each \(Q_I\) is \(\exists_I\) or
\(\forall_I\), which are quantifiers limited to \(I\).  A one-one map
\(\varphi\) of \(I\) onto \(I'\) preserves finiteness, so \(Q_I\) and
\(Q_{I'}\) are both existential or both general, and
\(Q_I(F)=Q_{I'}(F_\varphi)\) because \(\varphi\) maps \(F^{-1}(\vee)\) onto
\(F_\varphi^{-1}(\vee)\).

\emph{Transfer to all infinite sets.}  Let \(\mathrm Y\) contain the
quantifiers \(\exists,\forall,Q\) and be satisfiable in an infinite \(I\).
Let \(\mathrm Y^\circ\) replace every \(\mathbf Q\) by \(\exists\).  On an
infinite set \(Q\) and \(\exists\) coincide, so by induction on formulas
\(\mathrm Y\) and \(\mathrm Y^\circ\) have the same value at every
valuation on an infinite set.  Read the functional variables of
\(\mathrm Y^\circ\) as relation symbols of a finite signature and let
\(\sigma\) be the existential closure of \(\mathrm Y^\circ\) over its free
individual variables; a valuation on a set satisfying \(\mathrm Y^\circ\)
is the same as a structure on that set satisfying \(\sigma\), a
propositional function with support \(\{1,\ldots,k\}\) corresponding to a
\(k\)-ary relation.  \(\sigma\) has the infinite model given by \(I\), so
by \cref{thm:upward-ls} (axiom of choice) it has a model of every infinite
power; transporting one of power \(\overline J\) to \(J\) along a bijection
gives a \(J\)-valuation satisfying \(\mathrm Y^\circ\), hence
\(\mathrm Y\).

\emph{\(Q\) is not definable in terms of \(\exists,\forall\).}  Suppose
\(\mathrm D\), with the quantifiers \(\exists,\forall\) and the unique free
variable \(\mathrm F\), makes \((\mathbf Q\mathrm x)\mathrm F(\mathrm x)
\equiv\mathrm D\) true in every set.  Rename the bound variables of
\(\mathrm D\) away from \(\mathrm y\) and let
\[
 \mathrm Y=(\exists\mathrm y)\,\mathrm D[(\mathrm z=\mathrm y)/\mathrm F]
 \wedge(\exists\mathrm u)(\exists\mathrm v){\sim}(\mathrm u=\mathrm v).
\]
By \cref{thm:p03-w03-p03-042-val-invariance}(3), at a valuation giving
\(\mathrm y\) the value \(b\), the first conjunct's matrix has the value of
\(\mathrm D\) with \(\mathrm F\) interpreted as \(\{b\}\), which is
\(Q_I(\{b\})\).  On an infinite \(I\), \(Q_I(\{b\})=\exists_I(\{b\})=\vee\),
so \(\mathrm Y\) is true.  On a finite \(I\) with at least two elements,
\(Q_I(\{b\})=\forall_I(\{b\})=\wedge\) for every \(b\), so \(\mathrm Y\) is
false; on a one-element set the second conjunct is false.  \(\mathrm Y\) is
a sentence of first-order logic with identity and no non-logical symbols,
true in exactly the infinite sets.  Then every finite subset of
\(\{{\sim}\mathrm Y\}\cup\{\lambda_k:k\geq1\}\), where \(\lambda_k\) says
that there are at least \(k\) elements, has a finite model, so by
\cref{thm:fol-compactness} the whole set has a model; it is infinite and
satisfies \({\sim}\mathrm Y\), a contradiction.
\end{proof}
\end{annexitem}

\begin{annexitem}{annex:p03-w03-p03-042-monadic-identity}{thm:p03-w03-p03-042-monadic-identity}{What identity would change}
\begin{proof}
(1) Let \(I\) have two distinct elements \(a,b\), and let \(M\) be an
\(I\)-valuation with \([\mathrm F_h]_M=\varnothing\) for \(h\leq n\) and
\([\mathrm x_1]_M=a\); let \(M'\) agree with \(M\) except that
\([\mathrm x_2]_{M'}=b\), while \([\mathrm x_2]_M=a\).  Every atom
\(\mathrm F_h(\mathrm x_l)\) has the value \(\wedge\) at both, and each
closed formula has the same value at both, since by
\cref{thm:p03-w03-p03-042-val-invariance}(1) its value does not depend on
the individual variables and \(M,M'\) agree on all functional variables.
Hence \(\mathrm U\), a stroke combination of these, has the same value at
\(M\) and \(M'\), while \(\mathrm x_1=\mathrm x_2\) has the value \(\vee\)
at \(M\) and \(\wedge\) at \(M'\).  So \(\mathrm x_1=\mathrm x_2\equiv
\mathrm U\) is false at one of them.

(2) There are non-denumerably many sets of positive integers (Cantor's
theorem) and only denumerably many recursive ones (each is given by one of
denumerably many finite systems of recursion equations), so some set \(A\)
of positive integers is not recursive.  Let \(Q_I(X)=\vee\) iff \(X\) is
finite and \(\overline X\in A\), for every \(I\) and \(X\subseteq I\).  A
one-one map of \(I\) onto \(I'\) preserves the powers of subsets, so
\(Q\) is an unlimited quantifier.  For \(k\geq1\) let
\[
 \begin{aligned}
 \mathrm E_k={}&\Bigl[[\mathrm F_1(\mathrm x_1)\vee{\sim}\mathrm
 F_1(\mathrm x_1)]\wedge\bigwedge_{1\leq a<b\leq k}{\sim}(\mathrm
 x_a=\mathrm x_b)\Bigr]\\
 &\supset{\sim}(\mathbf Q^1\mathrm y)[\mathrm
 y=\mathrm x_1\vee\cdots\vee\mathrm y=\mathrm x_k],
 \end{aligned}
\]
with \(\mathrm y\) distinct from \(\mathrm x_1,\ldots,\mathrm x_k\); the
first conjunct of the antecedent is true at every valuation and only puts
\(\mathrm F_1\) into the formula, and for \(k=1\) the second conjunct is
absent.  At a valuation \(M\) on a set \(I\),
if \([\mathrm x_1]_M,\ldots,[\mathrm x_k]_M\) are distinct, the formula
\((\mathbf Q^1\mathrm y)[\ldots]\) has the value of \(Q_I\) at the
\(k\)-element set \(\{[\mathrm x_1]_M,\ldots,[\mathrm x_k]_M\}\), namely
\(\vee\) iff \(k\in A\).  Hence \(\mathrm E_k\) is true in a set \(I\)
with at least \(k\) elements iff \(k\notin A\), and it is true in every
set with fewer than \(k\) elements, where the antecedent is never
satisfied.  So for a denumerable \(I\), and also for truth in every set,
\(\mathrm E_k\) is true iff \(k\notin A\).  The map \(k\mapsto\mathrm
E_k\) is computable.  If either set of true formulas were recursive, the
complement of \(A\), and so \(A\), would be recursive.
\end{proof}
\end{annexitem}

\begin{annexitem}{annex:p03-w03-p03-042-mycielski}{thm:p03-w03-p03-042-mycielski}{Mycielski's lemma}
\begin{proof}
\emph{Injectivity.}  Let \(A_1\neq A_2\) be subsets of \(A_0\) containing
\(1\).  Exchanging the names if necessary, there is \(p\in A_1-A_2\).
As \(1\) lies in both sets, \(p\neq1\), so \(p=2^{3^n}\) for some
\(n\geq0\) and \(p\geq2\).  Then \(p\in H(A_1)\): take \(n_1=\cdots=
n_p=1\in A_1\), whose sum \(p\) lies in \(A_1\).

Suppose \(p\in H(A_2)\), with \(n_1,\ldots,n_p\in A_2\) and
\(n_0=n_1+\cdots+n_p\in A_2\); order them so that \(n_1\leq\cdots\leq
n_p\).  If all \(n_i\) were \(1\), then \(n_0=p\in A_2\), which is false.
So \(n_p>1\) and \(n_p=2^{3^s}\) with \(s\geq0\).  Since \(n_0\geq
(p-1)+n_p\geq p+1\geq3\), \(n_0\) is neither \(1\) nor \(2\), so
\(n_0=2^{3^t}\) with \(t\geq1\).  Also \(n_0\geq n_p+(p-1)>n_p\), so
\(3^t>3^s\) and \(s\leq t-1\).  From \(n_0\leq pn_p\),
\[
 p\geq\frac{n_0}{n_p}=2^{3^t-3^s}\geq2^{3^t-3^{t-1}}=2^{2\cdot3^{t-1}},
\]
and from \(n_0\geq p+1\), \(p<2^{3^t}\).  With \(p=2^{3^n}\) this gives
\(2\cdot3^{t-1}\leq3^n<3^t\).  The right inequality gives \(n\leq t-1\),
so \(3^n\leq3^{t-1}<2\cdot3^{t-1}\), contradicting the left one.  Hence
\(p\notin H(A_2)\) and \(H(A_1)\neq H(A_2)\).

\emph{A non-recursive \(H(A)\).}  The sets \(A\subseteq A_0\) with
\(1\in A\) correspond one-one to the subsets of the denumerable set
\(A_0-\{1\}\), so by Cantor's theorem they are not denumerable.  By the
first part, \(A\mapsto H(A)\) is one-one on them.  Each recursive set of
integers is given by one of denumerably many finite systems of recursion
equations, so there are only denumerably many recursive sets.  If every
\(H(A)\) were recursive, the map would inject a non-denumerable family
into a denumerable one.  Hence some \(H(A)\) is not recursive.  No choice
is used.
\end{proof}
\end{annexitem}

\begin{annexitem}{annex:p03-w03-p03-042-theorem-nine}{thm:p03-w03-p03-042-theorem-nine}{Mostowski's Theorem~9}
\begin{proof}
Let \(A\) be as in \cref{thm:p03-w03-p03-042-mycielski}, with \(H(A)\)
not recursive, and let \(Q_I(X)=\vee\) iff \(X\) is finite and
\(\overline X\in A\), for all \(I\) and \(X\subseteq I\).  As in
\cref{annex:p03-w03-p03-042-monadic-identity}, \(Q\) is an unlimited
quantifier.  Take \(s\geq3\), \(Q^1=\forall\), \(Q^2=\exists\),
\(Q^3=\cdots=Q^s=Q\).  Let \(I\) be denumerable and \(n\geq1\).

At an \(I\)-valuation \(M\), let \(X_k\) be the set of elements
satisfying \(\mathrm S_k\) (\(k\leq n\)).  These are cells of the
partition by constituents of order \(n\), hence pairwise disjoint, and
\(X_1\cup\cdots\cup X_n\) is the set satisfying \(\mathrm S_1\vee\cdots
\vee\mathrm S_n\).  By clause~4 of \cref{def:p03-w03-p03-042-calculus-s},
\(\mathit{val}_{MI}(\mathrm W_n)=\vee\) iff \(Q_I(X_k)=\vee\) for each
\(k\leq n\) and \(Q_I(X_1\cup\cdots\cup X_n)=\vee\).

If this holds, each \(X_k\) is finite with \(n_k=\overline{X_k}\in A\),
and the disjoint union has \(n_1+\cdots+n_n\) elements, which lies in
\(A\); so \(n\in H(A)\).  Conversely, let \(n\in H(A)\) with
\(n_1,\ldots,n_n\in A\) and \(n_1+\cdots+n_n\in A\).  As \(I\) is
infinite, it has pairwise disjoint subsets \(X_1,\ldots,X_n\) with
\(\overline{X_k}=n_k\).  Let \(e^{(k)}\in\{0,1\}^n\) be the index of
\(\mathrm S_k\), so \(e^{(1)},\ldots,e^{(n)}\) are the first \(n\) tuples
in lexicographic order; the tuple \((1,\ldots,1)\) is the last of the
\(2^n>n\) tuples, so it is none of them.  Put \([\mathrm F_h]_M=\bigcup\{
X_k:k\leq n,\ e^{(k)}_h=0\}\), and let the individual variables take
arbitrary values.  An element of \(X_k\) belongs to \([\mathrm F_h]_M\)
exactly when \(e^{(k)}_h=0\), so it satisfies \(\mathrm S_k\); an element
outside \(X_1\cup\cdots\cup X_n\) lies in no \([\mathrm F_h]_M\) and
satisfies the constituent with index \((1,\ldots,1)\).  So the set
satisfying \(\mathrm S_k\) is \(X_k\), and \(\mathit{val}_{MI}(\mathrm
W_n)=\vee\).  Thus \(\mathrm W_n\) is satisfiable in \(I\) iff \(n\in
H(A)\).

\(n\mapsto\mathrm W_n\) is computable.  If the set of formulas of
\((\mathrm S^M)\) satisfiable in \(I\) were recursive, \(H(A)\) would be
recursive.  By \cref{thm:p03-w03-p03-042-val-invariance}(2) the set does
not depend on which denumerable \(I\) is taken.  The formulas
\(\mathrm W_n\) contain no identity, so the same holds whether or not
identity is allowed in \((\mathrm S^M)\).
\end{proof}
\end{annexitem}

\begin{annexitem}{annex:p03-w03-p03-042-theorem-nine-true}{thm:p03-w03-p03-042-theorem-nine-true}{True formulas in a denumerable set}
\begin{proof}
\(\mathrm Z\) is satisfiable in \(I\) iff some \(I\)-valuation gives it
\(\vee\) iff \({\sim}\mathrm Z\), that is \(\mathrm Z|\mathrm Z\), is not
true in \(I\).  The map \(\mathrm Z\mapsto\mathrm Z|\mathrm Z\) is
computable and stays in \((\mathrm S^M)\).  If the set of formulas true in
a denumerable \(I\) were recursive, so would be the set of formulas
satisfiable in \(I\), against \cref{thm:p03-w03-p03-042-theorem-nine}.
\end{proof}
\end{annexitem}

\begin{annexitem}{annex:p03-w03-p03-042-ten-lemma-a}{thm:p03-w03-p03-042-ten-lemma-a}{Lemma~(a) of Theorem~10}
\begin{proof}
Identify \(F,G\) with the sets \(X,Y\subseteq I\) where they take the
value \(\vee\); then \(F\vee G\) is \(X\cup Y\), and \(Q^j_I(X)=\vee\) iff
\(\overline X\geq\mathfrak m_j\).  Since \(\mathfrak m_j\geq\mathfrak
m_1\geq1\), the empty set gets \(\wedge\).  If \(\overline X\geq
\mathfrak m_j\) then \(\overline{X\cup Y}\geq\overline X\geq\mathfrak
m_j\), and if \(\overline Y\geq\mathfrak m_j\) then \(\overline{X\cup
Y}\geq\overline Y\geq\mathfrak m_j\); so the right side \(\vee\) gives the left
side \(\vee\).

Conversely let \(\overline{X\cup Y}\geq\mathfrak m_j\).  If \(\mathfrak
m_j=1\), \(X\cup Y\) is nonempty, so \(X\) or \(Y\) is nonempty.  Let
\(\mathfrak m_j\) be infinite, and suppose \(\overline X<\mathfrak m_j\)
and \(\overline Y<\mathfrak m_j\).  By the axiom of choice the powers are
comparable; say \(\overline X\leq\overline Y\).  If \(Y\) were finite,
\(X\cup Y\) would be finite and could not have power at least an infinite
\(\mathfrak m_j\).  So \(Y\) is infinite, and, again by the axiom of
choice, \(\overline{X\cup Y}\leq\overline X+\overline Y\leq\overline Y+
\overline Y=\overline Y<\mathfrak m_j\), a contradiction.  So \(\overline
X\geq\mathfrak m_j\) or \(\overline Y\geq\mathfrak m_j\).
\end{proof}
\end{annexitem}

\begin{annexitem}{annex:p03-w03-p03-042-ten-lemma-b}{thm:p03-w03-p03-042-ten-lemma-b}{Lemma~(b) of Theorem~10}
\begin{proof}
Write L and R for the two sides, fix \(I\) and an \(I\)-valuation \(M\).
For \(y\in I\) let \(M(i,y)\) be as in
\cref{def:p03-w03-p03-042-calculus-s}.  Since \(\mathrm W_h\) does not
contain \(\mathrm x_i\), \(\mathit{val}_{M(i,y)I}(\mathrm W_h)=w_h:=
\mathit{val}_{MI}(\mathrm W_h)\) for every \(y\)
(\cref{thm:p03-w03-p03-042-val-invariance}(1)).  Let \(F_h\) be the set
of \(y\) with \(\mathit{val}_{M(i,y)I}(\mathrm Z_h)=\vee\).  The set of
\(y\) giving \((\mathrm Z_1\wedge\mathrm W_1)\vee\cdots\vee(\mathrm Z_n
\wedge\mathrm W_n)\) the value \(\vee\) is \(\bigcup\{F_h:w_h=\vee\}\).
By additivity, applied \(n-1\) times, and \(Q^j_I(\varnothing)=\wedge\)
when no \(w_h\) is \(\vee\),
\[
 \mathit{val}_{MI}(\mathrm L)=Q^j_I\Bigl(\bigcup_{w_h=\vee}F_h\Bigr)=
 \bigvee_{w_h=\vee}Q^j_I(F_h)=\bigvee_{h=1}^{n}[w_h\wedge Q^j_I(F_h)]
 =\mathit{val}_{MI}(\mathrm R),
\]
the last step because \(\mathit{val}_{MI}((\mathbf Q^j\mathrm x_i)
\mathrm Z_h)=Q^j_I(F_h)\).  As \(I,M\) were arbitrary, \(\mathrm L\equiv
\mathrm R\) is true.  This is Mostowski's computation, in which the step
\(w_h\wedge Q^j_I(F_h)=Q^j_I(F_h\wedge w_h)\) is the case distinction made
by the index set \(\{h:w_h=\vee\}\).
\end{proof}
\end{annexitem}

\begin{annexitem}{annex:p03-w03-p03-042-ten-lemma-c}{thm:p03-w03-p03-042-ten-lemma-c}{Lemma~(c) of Theorem~10}
Two facts about (S) are used, both immediate from the clauses of
\cref{def:p03-w03-p03-042-calculus-s}.  \emph{Replacement}: if
\(\mathrm Z_1\equiv\mathrm U_1\) is true, so are \((\mathbf Q^j\mathrm y)
\mathrm Z_1\equiv(\mathbf Q^j\mathrm y)\mathrm U_1\) (the two defining
sets in clause~4 coincide at every valuation) and, with \(\mathrm
Z_2\equiv\mathrm U_2\) true, \(\mathrm Z_1|\mathrm Z_2\equiv\mathrm
U_1|\mathrm U_2\).  \emph{Renaming}: for a constituent \(\mathrm C(\mathrm
y)\) and the fixed bound variable \(\mathrm x\) of
\cref{def:p03-w03-p03-042-ten-class}, \((\mathbf Q^j\mathrm y)\mathrm
C(\mathrm y)\equiv(\mathbf Q^j\mathrm x)\mathrm C(\mathrm x)\) is true,
because at every valuation both quantify the set of elements satisfying
the constituent, which does not depend on the values of the individual
variables.  Write \(\tau\) for \(\mathrm A|(\mathrm A|\mathrm A)\) with
\(\mathrm A=\mathrm A^1_{0\ldots0}\); it is in \(\mathfrak R\) and has the
value \({\sim}a\vee a=\vee\) everywhere, and \(\tau|\tau\) has the value
\(\wedge\) everywhere.
\begin{proof}
By induction on \(\mathrm Z\) we prove: for every finite list \(\vec v\)
of distinct individual variables containing the free variables of
\(\mathrm Z\), a formula \(\mathrm U\in\mathfrak R(\vec v)\) with
\(\mathrm Z\equiv\mathrm U\) true can be computed from \(\mathrm Z\) and
\(\vec v\).  Here \(\mathfrak R(\vec v)\) is
\cref{def:p03-w03-p03-042-ten-class} with \(\vec v\) for \(\mathrm
x_1,\ldots,\mathrm x_m\).  Mostowski's induction on the number of
quantifier symbols is this one with the atomic and stroke cases made
explicit.

\emph{Atoms.}  The atoms of \((\mathrm S^M)\) are \(\mathrm F_h(\mathrm
v)\) with \(\mathrm v\) in \(\vec v\); take \(\mathrm U=\mathrm Z\).

\emph{Stroke.}  If \(\mathrm U_1,\mathrm U_2\) serve for \(\mathrm
Z_1,\mathrm Z_2\), then \(\mathrm U_1|\mathrm U_2\) serves for
\(\mathrm Z_1|\mathrm Z_2\) by replacement.

\emph{Quantifier.}  Let \(\mathrm Z=(\mathbf Q^j\mathrm y)\mathrm Z_1\).
Let \(\vec v'\) be \(\vec v\) with \(\mathrm y\) deleted if it occurs.
The free variables of \(\mathrm Z_1\) are among \(\vec v',\mathrm y\);
the induction hypothesis gives \(\mathrm U_1\in\mathfrak R(\vec v',
\mathrm y)\) with \(\mathrm Z_1\equiv\mathrm U_1\) true, and by
replacement \(\mathrm Z\equiv(\mathbf Q^j\mathrm y)\mathrm U_1\) is true.
For \(e\in\{0,1\}^n\) let \(\mathrm U_1^e\) be \(\mathrm U_1\) with each
atom \(\mathrm F_h(\mathrm y)\) replaced by \(\tau\) if \(e_h=0\) and by
\(\tau|\tau\) if \(e_h=1\).  Then \(\mathrm U_1^e\in\mathfrak R(\vec v')\)
and it does not contain \(\mathrm y\).  Let \(\mathrm C_e(\mathrm y)\) be
the constituent with index \(e\).  The formula
\[
 \mathrm U_1\equiv\bigvee_{e\in\{0,1\}^n}[\mathrm C_e(\mathrm y)\wedge
 \mathrm U_1^e]
\]
is true: at a valuation, exactly one \(\mathrm C_e(\mathrm y)\) has the
value \(\vee\), and for that \(e\) every atom \(\mathrm F_h(\mathrm y)\)
has the value of its replacement in \(\mathrm U_1^e\), so \(\mathrm U_1\)
and \(\mathrm U_1^e\) have the same value (induction on the stroke
structure).  This is Mostowski's representation of \(\mathrm Z_1\) as a
sum of constituents in \(\mathrm y\) conjoined with members of
\(\mathfrak R(\vec v')\).  By replacement and
\cref{thm:p03-w03-p03-042-ten-lemma-b}, applied with the \(\mathrm
C_e(\mathrm y)\), which contain \(\mathrm y\), for the \(\mathrm Z_h\)
and the \(\mathrm U_1^e\), which do not, for the \(\mathrm W_h\) (Lemma~(b)
holds by \cref{thm:p03-w03-p03-042-ten-lemma-a}),
\[
 \mathrm Z\equiv\bigvee_{e}[\mathrm U_1^e\wedge(\mathbf Q^j\mathrm y)
 \mathrm C_e(\mathrm y)]
\]
is true, and by renaming and replacement so is \(\mathrm Z\equiv
\bigvee_e[\mathrm U_1^e\wedge\mathrm A^j_e]\).  The right side is in
\(\mathfrak R(\vec v')\subseteq\mathfrak R(\vec v)\), since \(\vee\) and
\(\wedge\) abbreviate stroke combinations.  Every step is a mechanical
rewriting of the given formula.
\end{proof}
\end{annexitem}

\begin{annexitem}{annex:p03-w03-p03-042-ten-lemma-d}{thm:p03-w03-p03-042-ten-lemma-d}{Lemma~(d) of Theorem~10, corrected}
For an \(I\)-valuation \(M\), let
\[
 \gamma_M(a^j_e)=\mathit{val}_{MI}(\mathrm A^j_e),\qquad
 \beta_M(b_{h,l})=\mathit{val}_{MI}(\mathrm F_h(\mathrm x_l)).
\]  Since \(\mathit{val}\) turns the stroke into
\({\sim}u\vee{\sim}v\), induction on the polynomial gives
\(\mathit{val}_{MI}(\mathrm U)=\Phi_{\mathrm U}(\gamma_M,\beta_M)\).
Write \(X_e\) for the set of elements satisfying the constituent with
index \(e\); then \(\gamma_M(a^j_e)=\vee\) iff \(\overline{X_e}\geq
\mathfrak m_j\).
\begin{proof}
\emph{If \(\mathrm U\) is not true, some admissible pair gives
\(\wedge\).}  Let \(\mathit{val}_{MI}(\mathrm U)=\wedge\); \(I\) is
nonempty since \(M\) exists.  \(\gamma_M\) is allowable: if
\(\overline{X_e}\geq\mathfrak m_{j+1}\) then \(\overline{X_e}\geq
\mathfrak m_j\).  If \(\mathfrak m_1=1\): \([\mathrm x_l]_M\) lies in
\(X_{c_l(\beta_M)}\), because \(c_l(\beta_M)\) records which
\([\mathrm F_h]_M\) contain it, so this set is nonempty and
\(\gamma_M(a^1_{c_l(\beta_M)})=\vee\); and some \(X_e\) is nonempty
because \(I\) is, so \(\gamma_M(a^1_e)=\vee\) for some \(e\).  So
\((\gamma_M,\beta_M)\) is admissible, and \(\Phi_{\mathrm U}(\gamma_M,
\beta_M)=\wedge\).

\emph{If an admissible pair gives \(\wedge\), \(\mathrm U\) is not
true.}  Let \((\gamma,\beta)\) be admissible with \(\Phi_{\mathrm U}
(\gamma,\beta)=\wedge\).  For each cell index \(e\) let \(j_e\) be the
largest \(j\leq s\) with \(\gamma(a^j_e)=\vee\), and \(j_e=0\) if there is
none.  By allowability, \(\gamma(a^j_e)=\vee\) iff \(1\leq j\leq j_e\).
Put \(\mu_e=\mathfrak m_{j_e}\) if \(j_e\geq1\); if \(j_e=0\) put
\(\mu_e=0\) when \(\mathfrak m_1=1\) and \(\mu_e=1\) when \(\mathfrak m_1\)
is infinite.  Take pairwise disjoint sets \(X_e\) of power \(\mu_e\) (for
instance \(\{e\}\times Y_e\) with \(\overline{Y_e}=\mu_e\)) and
\(I=\bigcup_eX_e\).

\(I\) is nonempty: if \(\mathfrak m_1\) is infinite every \(\mu_e\geq1\);
if \(\mathfrak m_1=1\), admissibility gives an \(e\) with
\(\gamma(a^1_e)=\vee\), so \(j_e\geq1\) and \(\mu_e\geq1\).  Likewise
\(X_{c_l(\beta)}\neq\varnothing\) for \(l\leq m\): for \(\mathfrak m_1\)
infinite as before, and for \(\mathfrak m_1=1\) because
\(\gamma(a^1_{c_l(\beta)})=\vee\).  Define \(M\) by \([\mathrm
F_h]_M=\bigcup\{X_e:e_h=0\}\), \([\mathrm x_l]_M\in X_{c_l(\beta)}\) for
\(l\leq m\), and arbitrary values in \(I\) for the other variables.  An
element of \(X_e\) lies in \([\mathrm F_h]_M\) iff \(e_h=0\), so it
satisfies exactly the constituent with index \(e\), and \(X_e\) is the set
satisfying that constituent.  Hence \(\mathit{val}_{MI}(\mathrm
F_h(\mathrm x_l))=\vee\) iff the \(h\)th entry of \(c_l(\beta)\) is \(0\)
iff \(\beta(b_{h,l})=\vee\); so \(\beta_M=\beta\).

For \(\gamma_M\): if \(j_e\geq1\), \(\overline{X_e}=\mathfrak m_{j_e}\),
and \(\mathfrak m_{j_e}\geq\mathfrak m_j\) iff \(j\leq j_e\), since the
\(\mathfrak m\)'s strictly increase (for \(j>j_e\),
\(\mathfrak m_{j_e}<\mathfrak m_j\) excludes \(\mathfrak m_{j_e}\geq
\mathfrak m_j\) by the Cantor--Bernstein theorem).  If \(j_e=0\),
\(\overline{X_e}=\mu_e<\mathfrak m_j\) for all \(j\geq1\): \(0<1\leq
\mathfrak m_j\) when \(\mathfrak m_1=1\), and \(1<\mathfrak m_1\leq
\mathfrak m_j\) when \(\mathfrak m_1\) is infinite.  In both cases
\(\gamma_M(a^j_e)=\vee\) iff \(j\leq j_e\) iff \(\gamma(a^j_e)=\vee\).
So \(\gamma_M=\gamma\) and \(\mathit{val}_{MI}(\mathrm U)=\Phi_{\mathrm
U}(\gamma,\beta)=\wedge\).

\emph{The special case.}  If \(\mathfrak m_1\) is infinite and \(m=0\),
admissible pairs are just allowable valuations, and the criterion is
property (T) of \(\Phi_{\mathrm U}\).  The counterexample in
\cref{def:p03-w03-p03-042-lemma-d-correction} shows that for
\(\mathfrak m_1=1\) the restriction to admissible pairs cannot be
dropped: the constant valuation \(\wedge\) is allowable but not
admissible.
\end{proof}
\end{annexitem}

\begin{annexitem}{annex:p03-w03-p03-042-theorem-ten}{thm:p03-w03-p03-042-theorem-ten}{Mostowski's Theorem~10}
\begin{proof}
Let \(\mathrm Z\) be a formula of \((\mathrm S^M)\), let \(\vec v\) list
its free individual variables, and let \(n\) be the largest index of a
functional variable in it (there is one, as every atom of \((\mathrm
S^M)\) contains one).  By \cref{thm:p03-w03-p03-042-ten-lemma-a} the
quantifiers satisfy the hypothesis of Lemmas~(b) and~(c), and
\cref{thm:p03-w03-p03-042-ten-lemma-c} computes \(\mathrm U\in\mathfrak
R(\vec v)\) with \(\mathrm Z\equiv\mathrm U\) true; then \(\mathrm Z\) is
true iff \(\mathrm U\) is, since the two have the same value at every
valuation.  By \cref{thm:p03-w03-p03-042-ten-lemma-d}, \(\mathrm U\) is
true iff \(\Phi_{\mathrm U}\) has the value \(\vee\) at each admissible
pair.  The pairs are the \(2^{s2^n+nm}\) assignments of truth values to
the variables \(a^j_e,b_{h,l}\) (\(m\) the length of \(\vec v\)), and
whether one is admissible depends only on \(s\), on whether
\(\mathfrak m_1=1\), and on the pair.  Evaluating \(\Phi_{\mathrm U}\) at
each admissible pair is a finite computation.  This procedure decides
truth; by the Church--Turing thesis (\cref{def:p04-church-turing}) the
set of true formulas of \((\mathrm S^M)\) is recursive.
\end{proof}
\end{annexitem}

\begin{annexitem}{annex:p03-w03-p03-042-eleven-lemma-a}{thm:p03-w03-p03-042-eleven-lemma-a}{Lemma~(a) of Theorem~11}
\begin{proof}
Fix \(I\) and \(M\).  \(\mathrm A\) and \(\mathrm B\) do not contain
\(\mathrm x\), so their values at \(M(\mathrm x,a)\) are their values
\(\alpha,\beta\) at \(M\) for every \(a\in I\)
(\cref{thm:p03-w03-p03-042-val-invariance}(1)).  Let \(X,Y\) be the sets
of \(a\) giving \(\mathrm M(\mathrm x),\mathrm N(\mathrm x)\) the value
\(\vee\).  The set defined by the matrix of the left side is
\(\varnothing\), \(Y\), \(X\) or \(X\cup Y\) according as \((\alpha,
\beta)\) is \((\wedge,\wedge)\), \((\wedge,\vee)\), \((\vee,\wedge)\) or
\((\vee,\vee)\).  On the right exactly one of the four guards
\({\sim}\mathrm A\wedge{\sim}\mathrm B,\ldots,\mathrm A\wedge\mathrm B\)
has the value \(\vee\), namely the one matching \((\alpha,\beta)\), and
the quantified formula next to it has the value \(Q^j_I\) of the same set:
for \(\mathrm A^j_0\) the set satisfying \(\mathrm F_1(\mathrm x)\wedge
{\sim}\mathrm F_1(\mathrm x)\) is \(\varnothing\).  So both sides have the
same value.  No property of \(Q^j\) is used.  Mostowski calls the proof
evident \cite[p.~32]{Mostowski1957}.
\end{proof}
\end{annexitem}

\begin{annexitem}{annex:p03-w03-p03-042-eleven-lemma-b}{thm:p03-w03-p03-042-eleven-lemma-b}{Lemma~(b) of Theorem~11}
\begin{proof}
Fix \(I\) and \(M\), and let \(\kappa_k=\mathit{val}_{MI}(\mathrm C_k)\),
which is also the value at every \(M(\mathrm x,a)\).  Let \(i^*\) be the
tuple with \(i^*_k=0\) iff \(\kappa_k=\vee\).  The guard \(\mathrm
C_1^{i_1}\wedge\cdots\wedge\mathrm C_p^{i_p}\) has the value \(\vee\) iff
\(i=i^*\), so the right side has the value of \((\mathbf Q^j\mathrm x)
\underset{i^*}{\Sigma}(\mathrm x)\).  For \(a\in I\), \(\mathrm Z\) has
the value \(\vee\) at \(M(\mathrm x,a)\) iff \(\mathrm M_0(\mathrm x)\)
does or some \(\mathrm N_k(\mathrm x)\) with \(\kappa_k=\vee\) does, that
is, iff \(\underset{i^*}{\Sigma}(\mathrm x)\) does.  The two quantified
formulas therefore apply \(Q^j_I\) to the same set and have the same
value.

\emph{Mostowski's route.}  For \(p=1\), Lemma~(a) with \(\mathrm A=
\mathrm D\vee{\sim}\mathrm D\), where Mostowski takes \(\mathrm D=
(\mathbf Q^j\mathrm x)\mathrm F_1(\mathrm x)\), \(\mathrm B=
\mathrm C_1\), \(\mathrm M=\mathrm M_0\), \(\mathrm N=\mathrm N_1\) gives
the claim, since \({\sim}\mathrm A\) is false at every valuation and
\(\mathrm A\wedge\mathrm M_0\) has the value of \(\mathrm M_0\).  For the
inductive step, Lemma~(a) with \(\mathrm A\) as before, \(\mathrm B=
\mathrm C_p\), \(\mathrm M(\mathrm x)=\mathrm M_0(\mathrm x)\vee
\bigvee_{k<p}[\mathrm C_k\wedge\mathrm N_k(\mathrm x)]\) and \(\mathrm
N=\mathrm N_p\) gives the true equivalence
\[
 (\mathbf Q^j\mathrm x)\mathrm Z\equiv\mathrm C_p^1\wedge(\mathbf
 Q^j\mathrm x)\mathrm M(\mathrm x)\vee\mathrm C_p^0\wedge(\mathbf
 Q^j\mathrm x)\{[\mathrm M_0(\mathrm x)\vee\mathrm N_p(\mathrm x)]\vee
 \textstyle\bigvee_{k<p}[\mathrm C_k\wedge\mathrm N_k(\mathrm x)]\}
\]
\cite[p.~33]{Mostowski1957}: as \(\mathrm A\) is true everywhere, the
left side of Lemma~(a) quantifies the same set as \((\mathbf Q^j\mathrm
x)\mathrm Z\), and \(\mathrm M\vee\mathrm N_p\) is satisfied by the same
elements as the formula in the last braces.  Then the induction hypothesis, applied to the two quantified
formulas with \(\mathrm M_0\) and with \(\mathrm M_0\vee\mathrm N_p\) as
the first disjunct, together with replacement, gives the claim for \(p\).
The direct argument above proves the same equivalence.
\end{proof}
\end{annexitem}

\begin{annexitem}{annex:p03-w03-p03-042-eleven-lemma-c}{thm:p03-w03-p03-042-eleven-lemma-c}{Lemma~(\(\mathrm c_n\)) of Theorem~11}
\begin{proof}
The induction of \cref{annex:p03-w03-p03-042-ten-lemma-c} is repeated,
with \(\tau=\mathrm A^1_0|(\mathrm A^1_0|\mathrm A^1_0)\); atoms and the
stroke are treated as there, and nothing about the quantifiers is used
until the quantifier step.  There \(\mathrm Z=(\mathbf Q^j\mathrm y)
\mathrm Z_1\), and as before \(\mathrm Z\equiv(\mathbf Q^j\mathrm y)
\mathrm Z'\) is true, where
\[
 \mathrm Z'=\mathrm M_0(\mathrm y)\vee\bigvee_{k=1}^{2^n}[\mathrm
 C_k\wedge\mathrm S_k(\mathrm y)],\qquad \mathrm M_0(\mathrm y)=\mathrm
 F_1(\mathrm y)\wedge{\sim}\mathrm F_1(\mathrm y),
\]
with \(\mathrm C_k=\mathrm U_1^{e}\in\mathfrak R(\vec v')\) for the index
\(e\) of \(\mathrm S_k\); adding the disjunct \(\mathrm M_0\), which is
false everywhere, does not change the value.  By
\cref{thm:p03-w03-p03-042-eleven-lemma-b} with \(p=2^n\) and \(\mathrm
N_k=\mathrm S_k\),
\[
 (\mathbf Q^j\mathrm y)\mathrm Z'\equiv\bigvee_{i\in\{0,1\}^p}\bigl[
 \mathrm C_1^{i_1}\wedge\cdots\wedge\mathrm C_p^{i_p}\wedge(\mathbf
 Q^j\mathrm y)\underset{i}{\Sigma}(\mathrm y)\bigr]
\]
is true.  Let \(K_i=\{k:i_k=0\}\).  If \(K_i\neq\varnothing\),
\(\underset{i}{\Sigma}(\mathrm y)\) and \(\bigvee_{k\in K_i}\mathrm
S_k(\mathrm y)\) are satisfied by the same elements at every valuation,
and the latter, with \(\mathrm x\) for \(\mathrm y\), is the canonical
representation of a non-vanishing polynomial \(\mathrm W\); so
\((\mathbf Q^j\mathrm y)\underset{i}{\Sigma}(\mathrm y)\equiv\mathrm
A^j_{\mathrm W}\) is true, the bound variable being immaterial as in the
renaming fact.  If \(K_i=\varnothing\), \(\underset{i}{\Sigma}=\mathrm
M_0\) is satisfied by no element and \((\mathbf Q^j\mathrm y)\mathrm
M_0(\mathrm y)\equiv\mathrm A^j_0\) is true.  Since \(\mathrm C_k^1=
\mathrm C_k|\mathrm C_k\in\mathfrak R(\vec v')\), replacement gives a
formula of \(\mathfrak R(\vec v')\subseteq\mathfrak R(\vec v)\)
equivalent to \(\mathrm Z\).  This is Mostowski's reduction to the two
forms of \(\mathrm Z_1\) \cite[p.~33]{Mostowski1957}.
\end{proof}
\end{annexitem}

\begin{annexitem}{annex:p03-w03-p03-042-eleven-lemma-d}{thm:p03-w03-p03-042-eleven-lemma-d}{Lemma~(d) of Theorem~11}
\begin{proof}
Let \(R_k\) be the set of \(x\in I\) satisfying \(\mathrm S_k\) at \(M\),
that is, giving \(\mathrm S_k(\mathrm x_i)\) the value \(\vee\) at
\(M(i,x)\), and \(\mathfrak m_k=\overline{R_k}\)
\cite[p.~34]{Mostowski1957}.  The \(R_k\) are disjoint with union \(I\),
so \(1^\circ\) holds.  If (15) are the canonical representations of
\(\mathrm W\) and \({\sim}\mathrm W\), the sets where \(\mathrm W\) is
true and false are \(R_{k_1}\cup\cdots\cup R_{k_p}\) and \(R_{l_1}\cup
\cdots\cup R_{l_q}\), of powers \(\sum\mathfrak m_{k_i}\) and
\(\sum\mathfrak m_{l_i}\) (the sets are disjoint).  By the definition of
\(Q_T\) (\cref{def:p03-w03-p03-042-qt}),
\(\mathit{val}_{MI}(\mathrm A^j_{\mathrm W})=T^j(\sum\mathfrak m_{k_i},
\sum\mathfrak m_{l_i})\), which is \(2^\circ\).  For \(\mathrm A^j_0\)
the defined set is empty and the value is \(T^j(0,\overline I)\), in
accordance with the convention.  For \(D_r\): if \(k\in D_r\), some
\([\mathrm x_l]_M\) lies in \(R_k\), so \(\mathfrak m_k\neq0\).
\end{proof}
\end{annexitem}

\begin{annexitem}{annex:p03-w03-p03-042-eleven-lemma-e}{thm:p03-w03-p03-042-eleven-lemma-e}{Lemma~(e) of Theorem~11}
\begin{proof}
Let \(\mathfrak m_1,\ldots,\mathfrak m_m\) satisfy \(1^\circ\) and
\(2^\circ\) for \(\gamma\) (and \(\mathfrak m_k\neq0\) for \(k\in D\) in
the second part).  Take pairwise disjoint sets of powers \(\mathfrak
m_k\); by \(1^\circ\) their union has the power of \(I\), and a one-one
map of the union onto \(I\) carries them to a partition \(I=R_1\cup
\cdots\cup R_m\) with \(\overline{R_k}=\mathfrak m_k\)
\cite[p.~34]{Mostowski1957}.  Let \([\mathrm F_i]_M\) be the union of the
\(R_k\) whose constituent \(\mathrm S_k\) has \(\mathrm F_i\) unnegated,
for \(i\leq n\).  Choose \([\mathrm x_l]_M\in R_{c(l)}\) for \(l\leq r\),
possible as \(\mathfrak m_{c(l)}\neq0\), and arbitrary elements of the
nonempty \(I\) for the other individual variables.  An element of
\(R_k\) is in \([\mathrm F_i]_M\) iff \(\mathrm F_i\) is unnegated in
\(\mathrm S_k\), so it satisfies \(\mathrm S_k\) and no other
constituent; this is display~(16).  For \(\mathrm W\) with the canonical
representations (15), the sets where \(\mathrm W\) is true and false are
\(R_{k_1}\cup\cdots\cup R_{k_p}\) and \(R_{l_1}\cup\cdots\cup R_{l_q}\),
so \(\mathit{val}_{MI}(\mathrm A^j_{\mathrm W})=T^j(\sum\mathfrak
m_{k_i},\sum\mathfrak m_{l_i})=\gamma(\mathrm A^j_{\mathrm W})\) by
\(2^\circ\); for \(\mathrm A^j_0\) both are \(T^j(0,\overline I)\).
\end{proof}
\end{annexitem}

\begin{annexitem}{annex:p03-w03-p03-042-theorem-eleven}{thm:p03-w03-p03-042-theorem-eleven}{Mostowski's Theorem~11}
For \(\mathrm U\in\mathfrak R(\mathrm x_1,\ldots,\mathrm x_r)\) let
\(\Phi_{\mathrm U}\) be its polynomial in variables for the
\(\mathrm A^j_0,\mathrm A^j_{\mathrm W}\) and \(b_{h,l}\) for \(\mathrm
F_h(\mathrm x_l)\), and for \(c:\{1,\ldots,r\}\to\{1,\ldots,m\}\) let
\(\beta_c(b_{h,l})=\vee\) iff \(\mathrm F_h\) is unnegated in \(\mathrm
S_{c(l)}\).  As in \cref{annex:p03-w03-p03-042-ten-lemma-d},
\(\mathit{val}_{MI}(\mathrm U)=\Phi_{\mathrm U}(\gamma_M,\beta_M)\), and
\(\beta_M=\beta_c\) when \([\mathrm x_l]_M\) satisfies \(\mathrm
S_{c(l)}\) for each \(l\).
\begin{proof}
\emph{Truth in a given \(I\).}  If \(I\) is empty, every formula is true
in it and the set is recursive.  Let \(I\) be nonempty.  Claim:
\(\mathrm U\in\mathfrak R(\mathrm x_1,\ldots,\mathrm x_r)\) is true in
\(I\) iff \(\Phi_{\mathrm U}(\gamma,\beta_c)=\vee\) for every
\((\gamma,D)\in\mathfrak K_I\) and every \(c:\{1,\ldots,r\}\to D\).  If
the condition fails at \((\gamma,D),c\), Lemma~(e)
(\cref{thm:p03-w03-p03-042-eleven-lemma-e}) gives \(M\) with
\(\gamma_M=\gamma\) and \(\beta_M=\beta_c\), so \(\mathrm U\) is false at
\(M\).  If \(\mathrm U\) is false at \(M\), let \(c(l)\) be the index of
the constituent satisfied by \([\mathrm x_l]_M\); Lemma~(d)
(\cref{thm:p03-w03-p03-042-eleven-lemma-d}) gives \((\gamma_M,D_r)\in
\mathfrak K_I\) with \(c\) mapping into \(D_r\), and \(\Phi_{\mathrm U}
(\gamma_M,\beta_c)=\wedge\).

\(\mathfrak K_I\) is a subset of the finite set of pairs of a function
from the \(s\cdot2^{2^n}\) formulas to truth values and a subset of
\(\{1,\ldots,2^n\}\).  For each such finite set \(\mathfrak K\) there is
an algorithm that, given \(\mathrm Z\) in \((\mathrm S^M_n)\), computes
\(\mathrm U\) by \cref{thm:p03-w03-p03-042-eleven-lemma-c} with \(\vec
v=\mathrm x_1,\ldots,\mathrm x_r\) listing the free variables, and tests
the condition above with \(\mathfrak K\) in place of \(\mathfrak K_I\).
Since \(\mathrm Z\equiv\mathrm U\) is true, the algorithm for
\(\mathfrak K=\mathfrak K_I\) decides truth in \(I\).  Hence the set of
formulas of \((\mathrm S^M_n)\) true in \(I\) is recursive (Church--Turing
thesis, \cref{def:p04-church-turing}).

\emph{Truth.}  \(\mathrm Z\) is true iff it is true in every nonempty
\(I\).  Let \(\mathfrak K=\bigcup_I\mathfrak K_I\) over nonempty \(I\), a
subset of the same finite set.  By the claim, \(\mathrm U\) is true iff
\(\Phi_{\mathrm U}(\gamma,\beta_c)=\vee\) for every \((\gamma,D)\in
\mathfrak K\) and \(c\) into \(D\); the same algorithm with this
\(\mathfrak K\) decides truth \cite[p.~35]{Mostowski1957}.

The argument shows that one of finitely many algorithms is correct
without identifying it: which pairs lie in \(\mathfrak K_I\) depends on
the values of the \(T^j\), about which nothing effective is assumed.
\end{proof}
\end{annexitem}

\begin{annexitem}{annex:p03-w03-p03-042-fermat}{thm:p03-w03-p03-042-fermat}{Mostowski's Fermat-prime example}
\begin{proof}
Let \(I\) be infinite and \(M\) an \(I\)-valuation, and let \(X\) be the
set of elements in \([\mathrm F]_M\).  The antecedent has the value
\(\vee\) iff \(\overline X\) is a prime of the form \(2^{2^n}+1\); then
\(X\) is finite.  By \cref{def:p03-w03-p03-042-examples}(b),
\((\overset{(k)}{\Sigma}\mathrm x)\mathrm F(\mathrm x)\) has the value
\(\vee\) iff \(\overline X=k\), so the consequent has the value \(\vee\)
iff \(\overline X\in\{3,5,17,257,65537\}\).  Hence the formula is false
at \(M\) iff \(\overline X\) is a prime of the form \(2^{2^n}+1\) outside
this list.  As \(I\) is infinite, every finite power is the power of some
\(X\subseteq I\), and each \(X\) is \([\mathrm F]_M\) for some \(M\).  So
the formula is true in \(I\) iff no prime of the form \(2^{2^n}+1\) lies
outside the list.  The five listed numbers are primes of that form, for
\(n=0,\ldots,4\); so the formula is true in \(I\) iff these are all such
primes.
\end{proof}
\end{annexitem}

\begin{annexitem}{annex:p03-w03-p03-042-mundici-pc-ec}{thm:p03-w03-p03-042-mundici-pc-ec}{Projective classes are elementary in \(\mathscr L^{\circledR}\)}
\begin{proof}
Following Mundici \cite[p.~213]{Mundici1985OtherQuantifiers}.  Let \(K\)
be the class of \(\tau\)-reducts of the models of a first-order sentence
\(\theta_0\) of \(\tau'\supseteq\tau\); only the finitely many symbols of
\(\tau'-\tau\) occurring in \(\theta_0\) matter, so \(\tau'-\tau\) is
finite.  Throughout, ``\(\mathfrak A\models\exists\bar h\,\beta\)'' for
new function symbols \(\bar h\) means that some expansion of \(\mathfrak
A\) by functions \(\bar h\) of the right arities satisfies \(\beta\).

\emph{Step 1: relations to functions.}  Let \(\delta(x)\) be the
disjunction of the conjunctions \(\bigwedge\Delta\), one for each
one-element \(\tau\)-structure in \(K\), where \(\Delta\) is the set of the
formulas \(R(x,\ldots,x)\) and \(\neg R(x,\ldots,x)\), \(R\) a relation
symbol of \(\tau\) occurring in \(\theta_0\), true of its element (in a
one-element structure every term denotes that element, so these atoms fix
the structure as far as \(\theta_0\) can see it); \(\delta\) is a finite
disjunction, and a one-element
\(\tau\)-structure is in \(K\) iff it satisfies \(\exists x\,\delta(x)\),
since \(K\) is closed under isomorphism and membership depends only on the
symbols of \(\theta_0\).  For each relation symbol \(R\in\tau'-\tau\) of
arity \(k\) take a new \(k\)-ary function symbol \(g_R\) (a constant if
\(k=0\)) and new constants \(c_0,c_1\); let \(\theta_0^\ast\) replace
each atom \(R(\bar t)\) by \(g_R(\bar t)=c_1\), and keep the function
symbols of \(\tau'-\tau\).  Put
\[
 \theta=\bigl(\exists x\forall y\,x=y\wedge\exists x\,\delta(x)\bigr)\vee
 \bigl(c_0\neq c_1\wedge\theta_0^\ast\bigr).
\]
Then \(K=\{\mathfrak A\mid\mathfrak A\models\exists\bar g\,\theta\}\), with
\(\bar g\) the new function and constant symbols together with those of
\(\tau'-\tau\): on a one-element structure the second disjunct fails and
the first is membership in \(K\); on a structure with at least two
elements, an expansion satisfying \(\theta_0\) gives one satisfying
\(\theta\) (choose \(c_0\neq c_1\) and \(g_R(\bar a)=c_1\) iff
\(R(\bar a)\)), and conversely an expansion satisfying the second disjunct
gives \(R(\bar a)\) iff \(g_R(\bar a)=c_1\), an expansion satisfying
\(\theta_0\).

\emph{Step 2: Skolem form.}  Put \(\theta\) in prenex form.  By the axiom
of choice, an expansion satisfies the prenex \(\forall\bar y_1\exists
z_1\forall\bar y_2\exists z_2\cdots\alpha_0\) iff it has a further
expansion by Skolem functions \(f_j\), each a function of the universally
quantified variables before \(z_j\), satisfying \(\forall\bar y\,
\alpha_0(\ldots,f_j(\bar y_1\cdots\bar y_j),\ldots)\).  So
\(K=\{\mathfrak A\mid\mathfrak A\models\exists h_1\cdots h_p\,\forall
x_1\cdots x_m\,\alpha\}\) with \(\alpha\) quantifier-free; the \(h_i\) are
the \(\bar g\) and the Skolem functions.

\emph{Step 3: constants.}  If some \(h_i\) is a constant, add a universal
variable \(w\) if \(m=0\), replace each occurrence of \(h_i\) by \(h_i'(x_1)\)
(or \(h_i'(w)\)) for a new unary \(h_i'\), and conjoin \(h_i'(u)=h_i'(u')\)
with new universal variables \(u,u'\).  Any value of the constant gives a
constant \(h_i'\), and a constant \(h_i'\) gives the constant.  Now every
\(h_i\) has arity at least \(1\).

\emph{Step 4: arguments are variables.}  If an \(h\)-term (a term with
outermost symbol among the \(h_i\)) has an argument that is not a variable,
pick such an argument \(s\) of minimal size and replace \(\forall\bar x\,
\beta(s)\) by \(\forall\bar x\,\forall t\,[t=s\to\beta(t)]\), \(t\) new,
replacing that one occurrence of \(s\): both say the same in every
expansion, as Mundici's instance \(\forall y,z\,\beta(f(g(y,z),h(y)))\)
and \(\forall y,z,t,u[t=g(y,z)\wedge u=h(y)\to\beta(f(t,u))]\) shows
\cite[p.~213]{Mundici1985OtherQuantifiers}.  Each replacement lowers the
number of function-symbol occurrences inside arguments of \(h\)-terms, so
the process stops; afterwards every \(h\)-term is \(h_i(\bar v)\) with
variables \(\bar v\).  Function symbols of \(\tau\) may still be applied
to \(h\)-terms.

\emph{Step 5: distinct, unshared argument variables.}  For each
occurrence of a variable \(v\) as an argument of an \(h\)-term that is not
the first argument occurrence of \(v\) in an \(h\)-term, replace that
occurrence by a new variable \(v'\) and replace \(\forall\bar x\,\beta\) by
\(\forall\bar x\,\forall v'[v=v'\to\beta']\).  Afterwards the argument
variables of each \(h\)-term occurrence are pairwise distinct and occur in
no other \(h\)-term occurrence.

\emph{Step 6: one term per function symbol.}  If \(h\) occurs as
\(h(\bar u^1),\ldots,h(\bar u^q)\) with \(q\geq2\), take new
\(h^1,\ldots,h^q\), replace the \(l\)th occurrence by \(h^l(\bar u^l)\), and
conjoin \(\bigwedge_{l<l'}(\bar u^l=\bar u^{l'}\to h^l(\bar u^l)=
h^{l'}(\bar u^{l'}))\), which is Mundici's device in his example
\cite[p.~213]{Mundici1985OtherQuantifiers}.  Given \(h\), take all \(h^l=
h\).  Conversely, the tuples \(\bar u^l\) are distinct variables
universally quantified independently, so the conjunct makes all \(h^l\)
the same function, which serves as \(h\).

\emph{Step 7: reading off the prefix.}  Now \(K=\{\mathfrak A\mid\mathfrak
A\models\exists h_1\cdots h_r\forall\bar x\,\alpha\}\), each \(h_i\)
occurring only in the single term \(h_i(\bar u^i)\), \(\bar u^i\) a tuple of
\(n_i\geq1\) distinct variables, disjoint from the others.  Let \(\chi\)
replace \(h_i(\bar u^i)\) by a new variable \(y^i\), and let \(Q\) have the
rows \(\forall\bar u^i\,\exists y^i\) (\(i\le r\)) and a row \(\forall
w\,\exists y_w\) with \(y_w\) not in \(\chi\) for each remaining variable
\(w\) of \(\bar x\) (and one such row if there is no row at all).  \(Q\) is
a partially ordered quantifier with pattern of pairs \(\langle
n_i,1\rangle\) and \(\langle1,1\rangle\), and \(\chi\in\mathscr
L_{\omega\omega}\) is quantifier-free.  By
\cref{def:p03-w03-p03-042-mundici-po-logic}, \(\mathfrak A\models Q\chi\)
iff there are \(f_i\) with \(\chi\) true whenever \(y^i=f_i(\bar u^i)\);
for the dummy rows any functions do; so \(\mathfrak A\models Q\chi\) iff
\(\mathfrak A\models\exists\bar h\,\forall\bar x\,\alpha\).  Hence
\(K=\mathrm{Mod}_{\mathscr L^{\circledR}}(Q\chi)\).
\end{proof}
\end{annexitem}

\begin{annexitem}{annex:p03-w03-p03-042-mundici-q-zero}{thm:p03-w03-p03-042-mundici-q-zero}{\(Q_0\) in \(\mathscr L(Q^{\mathrm H})\)}
\begin{proof}
Mundici states the equivalence without proof
\cite[p.~216]{Mundici1985OtherQuantifiers}.  Let \(X\) be the set
defined by \(\varphi\).

From right to left: \(u=v\) gives \(f(u)=g(u)\), so \(f=g\), and then
\(f(u)=f(v)\) implies \(u=v\): \(f\) is injective.  \(f\) maps \(X\) into
\(X-\{t\}\) with \(t\in X\), so \(t,f(t),f(f(t)),\ldots\) are pairwise
distinct (if \(f^i(t)=f^j(t)\) with \(i<j\), injectivity gives
\(t=f^{j-i}(t)\), which is not \(t\)), and \(X\) is infinite.

From left to right: let \(X\) be infinite.  By the axiom of choice take
distinct \(a_0,a_1,\ldots\) in \(X\); put \(t=a_0\), \(f(a_k)=a_{k+1}\),
\(f(u)=u\) for other \(u\), and \(g=f\).  \(f\) is injective, so
\(u=v\leftrightarrow f(u)=g(v)\); for \(u\in X\), \(f(u)\in X\) and
\(f(u)\neq a_0\).

The matrix written in the Henkin prefix uses \(y=f(u)\), \(y'=g(v)\), and
its three uses of \(f(u)\) are all \(y\); so by
\cref{def:p03-w03-p03-042-mundici-po-logic} the right side is an
\(\mathscr L(Q^{\mathrm H})\)-formula, and \(Q_0\) is definable in
\(\mathscr L(Q^{\mathrm H})\).
\end{proof}
\end{annexitem}

\begin{annexitem}{annex:p03-w03-p03-042-mundici-walkoe}{thm:p03-w03-p03-042-mundici-walkoe}{Walkoe's observation for the quantifiers \(Q_\pi\)}
\begin{proof}
Mundici reports the observation without proof
\cite[p.~213]{Mundici1985OtherQuantifiers}.  All translations below go by
induction on formulas: atomic formulas are kept, the translation commutes
with the connectives and \(\exists\), and only the quantifier case needs
an argument.  If in that case the new formula has the same free variables
as the old one and is equivalent to it in every structure under every
assignment (given that the translated matrix is), the translation
\(\varphi\mapsto\varphi^\ast\) satisfies
\(\mathfrak A\models\varphi[s]\iff\mathfrak A\models\varphi^\ast[s]\) for
all \(\mathfrak A,s\), which gives \(\leq\) of
\cref{def:p03-w03-p03-042-mundici-comparison}.

(1) Let \(\pi=\{\langle n,m\rangle\}\).  By
\cref{def:p03-w03-p03-042-mundici-po-logic}, \(\mathfrak A\models
Q_\pi\chi[s]\) iff there is \(f:A^n\to A^m\) such that \(\chi\) holds
whenever \(\bar y=f(\bar x)\).  If such an \(f\) exists, then
\(\forall\bar x\exists\bar y\,\chi\) holds, with \(\bar y=f(\bar x)\).
Conversely, if \(\forall\bar x\exists\bar y\,\chi\) holds, the axiom of
choice gives a function choosing for each \(\bar x\in A^n\) one witness
\(\bar y\in A^m\).  So \(Q_\pi\chi\) may be replaced by
\(\forall\bar x\exists\bar y\,\chi\), and \(\mathscr
L(Q_\pi)\leq\mathscr L_{\omega\omega}\); the converse inclusion is
trivial.

(2) Let \(\pi=\{\langle n_1,m_1\rangle,\ldots,\langle n_r,m_r\rangle\}\)
with \(r\geq2\).  Translate \(Q^{\mathrm H}x\,x'\,y\,y'\,\psi\), with
\(\psi^\ast\) the translation of \(\psi\), into \(Q_\pi\psi^\ast\) where
the prefix \(Q_\pi\) has \(x\) as the first universally and \(y\) as the
first existentially quantified variable of row~\(1\), \(x'\) and \(y'\) in
the same places in row~\(2\), and new variables, not occurring in
\(\psi^\ast\), everywhere else.  Fix \(\mathfrak A\) and \(s\); \(A\neq
\varnothing\), and let \(a\in A\).

If \(f,g:A\to A\) witness \(Q^{\mathrm H}\), put
\(f_1(x,\bar w)=(f(x),a,\ldots,a)\), \(f_2(x',\bar w)=(g(x'),a,\ldots,a)\)
and let \(f_i\) (\(i\geq3\)) be constant with value \((a,\ldots,a)\).
Whenever every \(\bar y^i\) takes the value \(f_i(\bar x^i)\), the
variables \(y,y'\) take the values \(f(x),g(x')\), so \(\psi^\ast\) holds,
its other variables not occurring in it.  Conversely, if \(f_1,\ldots,f_r\)
witness \(Q_\pi\psi^\ast\), let \(f(b)\) be the first coordinate of
\(f_1(b,a,\ldots,a)\) and \(g(b)\) the first coordinate of
\(f_2(b,a,\ldots,a)\).  Given \(x,x'\), assign \(a\) to the other
universally quantified variables of rows~\(1\) and~\(2\), anything to
those of the other rows, and let each \(\bar y^i\) be \(f_i(\bar x^i)\);
then \(y=f(x)\), \(y'=g(x')\), and \(\psi^\ast\) holds.  So \(\mathscr
L(Q^{\mathrm H})\leq\mathscr L(Q_\pi)\).

\(\mathscr L(Q^{\mathrm H})\) is not \((\omega,\omega)\)-compact
(\cref{thm:p03-w03-p03-042-mundici-not-compact}, whose proof does not use
the present theorem), so neither is any \(\mathscr L'\geq\mathscr
L(Q^{\mathrm H})\): a countable set of \(\mathscr L(Q^{\mathrm
H})\)-sentences with finitely satisfiable but unsatisfiable union
translates sentence by sentence into such a set of \(\mathscr
L'\)-sentences with the same models.  First-order logic is
\((\omega,\omega)\)-compact (\cref{thm:fol-compactness}), so if
\(\mathscr L(Q_\pi)\equiv\mathscr L_{\omega\omega}\) then \(\mathscr
L(Q^{\mathrm H})\leq\mathscr L_{\omega\omega}\), which is impossible.

The final sentence of the theorem combines (1) and (2).
\end{proof}
\end{annexitem}

\begin{annexitem}{annex:p03-w03-p03-042-mundici-pairing}{thm:p03-w03-p03-042-mundici-pairing}{Contracting partially ordered prefixes with a pairing function}
\begin{proof}
Following Mundici \cite[p.~214]{Mundici1985OtherQuantifiers}.  Let
\(\mathfrak M\) be a \(\tau\)-structure whose \(\tau_0\)-reduct is a model
of \(T\), and \(J:M^2\to M\) the bijection whose graph \(\rho\) defines.
Since \(M\neq\varnothing\), fix \(a\in M\) where needed.

\emph{Tuple codes.}  For \(k\geq1\) let \(\rho_1(x_1,z)\) be \(x_1=z\) and
\(\rho_{k+1}(x_1,\ldots,x_{k+1},z)\) be \(\exists w\,[\rho_k(x_2,\ldots,
x_{k+1},w)\wedge\rho(x_1,w,z)]\), with \(w\) new.  By induction on \(k\),
\(\rho_k\) defines in \(\mathfrak M\) the graph of a bijection
\(c_k:M^k\to M\), with \(c_1\) the identity and
\(c_{k+1}(b_1,\ldots,b_{k+1})=J(b_1,c_k(b_2,\ldots,b_{k+1}))\): a
composition of the bijection \(\mathrm{id}\times c_k:M\times M^k\to M^2\)
with \(J\).  The formulas \(\rho_k\) are first-order and do not depend
on \(\mathfrak M\).

The translation goes by induction on formulas as in the proof of
\cref{thm:p03-w03-p03-042-mundici-walkoe}, now with equivalence required
only in the structures \(\mathfrak M\) above.  Let \(\varphi=Q\chi\) with
\(Q\) of pattern \(\pi=\{\langle n_1,m_1\rangle,\ldots,\langle
n_r,m_r\rangle\}\), rows \(\forall\bar x^i\exists\bar y^i\), and let
\(\chi^\ast\in\mathscr L(Q^{\mathrm H})\) be the translation of \(\chi\).
If \(r=1\), replace \(\varphi\) by \(\forall\bar x^1\exists\bar
y^1\,\chi^\ast\), as in part~(1) of
\cref{thm:p03-w03-p03-042-mundici-walkoe} (the axiom of choice).  Let
\(r\geq2\).

\emph{Step A: one variable per row.}  For any \(\theta\) and any prefix
\(Q\) with rows \(\forall\bar x^i\exists\bar y^i\), \(i\le r\), take new
variables \(u_1,v_1,\ldots,u_r,v_r\) and put
\begin{multline*}
 \theta_A=\forall\bar x^1\bar y^1\cdots\bar x^r\bar y^r\\
 \Bigl[\bigwedge_{i\le r}\bigl(\rho_{n_i}(\bar x^i,u_i)\wedge
 \rho_{m_i}(\bar y^i,v_i)\bigr)\to\theta\Bigr],
\end{multline*}
and let \(Q_A\) have the rows \(\forall u_i\exists v_i\), \(i\le r\).
Claim: \(\mathfrak M\models Q\theta\iff\mathfrak M\models Q_A\theta_A\)
under every assignment.  If \(f_i:M^{n_i}\to M^{m_i}\) witness \(Q\theta\),
let \(g_i=c_{m_i}\circ f_i\circ c_{n_i}^{-1}\).  For any values of the
\(u_i\), with \(v_i=g_i(u_i)\), the antecedent of \(\theta_A\) forces
\(\bar x^i=c_{n_i}^{-1}(u_i)\) and \(\bar y^i=c_{m_i}^{-1}(v_i)=f_i(\bar
x^i)\), so \(\theta\) holds.  If \(g_i\) witness \(Q_A\theta_A\), let
\(f_i=c_{m_i}^{-1}\circ g_i\circ c_{n_i}\).  For any \(\bar x^i\), put
\(u_i=c_{n_i}(\bar x^i)\), \(v_i=g_i(u_i)\), \(\bar y^i=f_i(\bar x^i)\);
then the antecedent of \(\theta_A\) holds, so \(\theta\) holds.

\emph{Step B: two rows.}  Let \(Q_A\) have rows \(\forall u_i\exists
v_i\), \(i\le r\), and put
\[
 \theta_B=\theta\wedge\bigwedge_{i\le r}(u_i=u'_i\to v_i=v'_i),
\]
with new \(u'_i,v'_i\), and let \(Q_B\) have the two rows
\[
 \forall u_1\cdots u_r\,\exists v_1\cdots v_r\quad\text{and}\quad
 \forall u'_1\cdots u'_r\,\exists v'_1\cdots v'_r,
\]
which is Mundici's equivalence
\cite[p.~214]{Mundici1985OtherQuantifiers}.  Claim: \(Q_A\theta\) and
\(Q_B\theta_B\) are equivalent in every structure.  If \(g_1,\ldots,g_r\)
witness \(Q_A\theta\), let \(F(\bar u)=F'(\bar u)=(g_1(u_1),\ldots,
g_r(u_r))\); then \(\theta\) holds at \(\bar v=F(\bar u)\), and
\(u_i=u'_i\) gives \(v_i=g_i(u_i)=g_i(u'_i)=v'_i\).  Conversely, let
\(F,F':M^r\to M^r\) witness \(Q_B\theta_B\), and write \(F_i\) for the
\(i\)th coordinate.  For all \(\bar u,\bar u'\) with \(u_i=u'_i\),
\(F_i(\bar u)=F'_i(\bar u')\); taking \(\bar u'=\bar u\) gives
\(F_i(\bar u)=F'_i(\bar u)\), so \(F_i(\bar u)=F_i(\bar u')\) whenever
\(u_i=u'_i\).  Hence \(g_i(b)=F_i(b,\ldots,b)\) satisfies
\(g_i(u_i)=F_i(\bar u)\) for every \(\bar u\), and \(\theta\) holds at
\(\bar v=(g_1(u_1),\ldots,g_r(u_r))\), as \(\theta_B\) holds there.

\emph{Step C: the Henkin prefix.}  \(Q_B\) has pattern
\(\{\langle r,r\rangle,\langle r,r\rangle\}\).  Applying Step~A to
\(Q_B\theta_B\) gives \(Q_{BA}(\theta_B)_A\) with two rows
\(\forall u\exists v\), \(\forall u'\exists v'\), which is
\(Q^{\mathrm H}u\,u'\,v\,v'\,(\theta_B)_A\).

So \(\varphi\) is replaced by \(Q^{\mathrm H}\) applied to
\(((\chi^\ast_A)_B)_A\), where the first Step~A is applied to
\(Q\chi^\ast\).  This formula is in \(\mathscr L(Q^{\mathrm H})(\tau)\):
it is built from \(\chi^\ast\), from \(\rho\), with \(\rho\) a
\(\tau_0\)-formula, by first-order operations and one \(Q^{\mathrm H}\).
Its free variables are those of \(\varphi\), since all variables
introduced are bound, and by the three claims it is equivalent to
\(\varphi\) in \(\mathfrak M\) under every assignment.
\end{proof}
\end{annexitem}

\begin{annexitem}{annex:p03-w03-p03-042-mundici-peano}{thm:p03-w03-p03-042-mundici-peano}{\(\mathscr L^{\circledR}\) and \(\mathscr L(Q^{\mathrm H})\) over Peano arithmetic}
\begin{proof}
By \cref{lem:p05-pa-pair-beta}, \(\PA\) proves that the pairing graph
\(\mathsf P(a,b,w)\) is total and single-valued in \((a,b)\) and that each
\(w\) has unique components \(a,b\) with \(\mathsf P(a,b,w)\).  These
\(\PA\)-theorems hold in every model \(\mathfrak M\) of \(\PA\), so
\(\mathsf P\) defines in \(\mathfrak M\) the graph of a function from
\(M^2\) to \(M\) that is onto (existence of components) and one-one
(uniqueness of components).  Apply
\cref{thm:p03-w03-p03-042-mundici-pairing} with \(T=\PA\),
\(\tau_0=\Lang_A\) and \(\rho=\mathsf P\) to a sentence \(\varphi\): the
sentence \(\psi\) obtained is true in exactly the same models of \(\PA\).
\end{proof}
\end{annexitem}

\begin{annexitem}{annex:p03-w03-p03-042-mundici-not-compact}{thm:p03-w03-p03-042-mundici-not-compact}{\(\mathscr L(Q^{\mathrm H})\) is neither \((\omega,\omega)\)-compact nor axiomatizable}
\begin{proof}
Mundici's proof is one paragraph \cite[p.~216]{Mundici1985OtherQuantifiers};
the details follow.  Write \(\bar n\) for the numeral \(S^n0\).

(1) By \cref{thm:p03-w03-p03-042-mundici-q-zero}, \(Q_0y\,\exists z\,
(y+z=x)\) is equivalent to an \(\mathscr L(Q^{\mathrm H})\)-formula with
the parameter \(x\), so \(\sigma_{\mathbb N}\in\mathscr L(Q^{\mathrm
H})(\Lang_A)\) and \(\mathfrak M\models\sigma_{\mathbb N}\) iff
\(\mathfrak M\) satisfies the axioms of \(\mathsf Q\), commutativity of
\(+\), and, for every \(b\in M\), the set \(D_b=\{c\in M\mid c+d=b\text{
for some }d\in M\}\) is finite.  In \(\mathcal N\) the axioms are true,
and \(D_b=\{0,\ldots,b\}\) is finite.  So \(\mathcal N\) and every
structure isomorphic to it are models.

Conversely let \(\mathfrak M\models\sigma_{\mathbb N}\) and
\(h(n)=\bar n^{\mathfrak M}\).  (a) \(h(0)=0^{\mathfrak M}\) and
\(h(n+1)=S^{\mathfrak M}h(n)\).  (b) \(h(n+m)=h(n)+^{\mathfrak M}h(m)\) by
induction on \(m\), using \(x+0=x\) and \(x+Sy=S(x+y)\); then \(h(nm)=
h(n)\times^{\mathfrak M}h(m)\) by induction on \(m\), using \(x\times0=0\)
and \(x\times Sy=(x\times y)+x\) together with (b).  (c) \(h\) is
one-one: if \(n<m\) and \(h(n)=h(m)\), cancelling \(S\) \(n\) times (the
axiom \(Sx=Sy\to x=y\)) gives \(0^{\mathfrak M}=h(m-n)=S^{\mathfrak
M}h(m-n-1)\), against \(Sx\neq0\).  (d) \(h\) is onto.  Suppose \(b\in M\)
is not a value of \(h\).  By induction on \(n\), \(h(n)+d_n=b\) for some
\(d_n\in M\): for \(n=0\), \(0+b=b+0=b\) by commutativity and \(x+0=x\);
given \(h(n)+d_n=b\), \(d_n\neq0\), since otherwise \(b=h(n)+0=h(n)\), so
\(d_n=Sw\) for some \(w\) by the axiom \(x\neq0\to\exists y\,x=Sy\), and
\begin{align*}
 h(n+1)+w&=Sh(n)+w=w+Sh(n)=S(w+h(n))\\
 &=S(h(n)+w)=h(n)+Sw=b .
\end{align*}
So every \(h(n)\) lies in \(D_b\), and \(D_b\) is infinite by (c),
contradicting \(\sigma_{\mathbb N}\).  By (a)--(d), \(h\) is an
isomorphism of \(\mathcal N\) onto \(\mathfrak M\).

(2) Let \(c\) be a new constant and \(\Sigma=\{\sigma_{\mathbb N}\}\cup
\{c\neq\bar n\mid n\in\mathbb N\}\), a countable set of
\(\mathscr L(Q^{\mathrm H})(\Lang_A\cup\{c\})\)-sentences.  A finite
subset mentions only \(c\neq\bar n\) for \(n<k\), and \(\mathcal N\) with
\(c\) interpreted as \(k\) is a model.  If \(\mathfrak M\models\Sigma\),
its \(\Lang_A\)-reduct satisfies \(\sigma_{\mathbb N}\), whose truth
depends only on that reduct, so by (1) there is an isomorphism \(h\) as
above, and \(c^{\mathfrak M}=h(n)=\bar n^{\mathfrak M}\) for some \(n\),
against \(c\neq\bar n\).

(3) For a first-order \(\Lang_A\)-sentence \(\theta\), \(\sigma_{\mathbb
N}\to\theta\) is valid iff \(\theta\) holds in every model of
\(\sigma_{\mathbb N}\) iff \(\mathcal N\models\theta\), by (1) and
isomorphism invariance.  If the valid sentences of \(\mathscr
L(Q^{\mathrm H})(\Lang_A)\) had a recursively enumerable set of codes,
then, the passage from \(\ulcorner\theta\urcorner\) to the code of
\(\sigma_{\mathbb N}\to\theta\) being computable
(\cref{def:p03-w03-p03-042-mundici-comparison}), the set of codes of true
\(\Lang_A\)-sentences would be recursively enumerable.  By the last clause
of \cref{thm:representability} it would then be defined in \(\mathcal N\)
by a \(\Sigma^0_1\) formula, which would be an arithmetical truth
definition, contradicting \cref{thm:tarski-undefinability}.
\end{proof}
\end{annexitem}

\begin{annexitem}{annex:p03-w03-p03-042-mundici-pure-identity}{thm:p03-w03-p03-042-mundici-pure-identity}{Pure identity formulas on infinite sets}
\begin{proof}
Mundici gives the induction and names its tools
\cite[p.~216]{Mundici1985OtherQuantifiers}; the details follow.  By
induction on \(\varphi\): for atomic \(\varphi\) put
\(\hat\varphi=\varphi\), and let \(\hat{}\) commute with the connectives
and \(\exists\); the equivalence on infinite sets passes through these
steps.

Let \(\varphi(\bar z)=Q^{\mathrm H}x\,x'\,y\,y'\,\psi\), with
\(\bar z=(z_1,\ldots,z_k)\) the free variables of \(\varphi\), and let
\(\hat\psi\) be given for \(\psi\).  On an infinite set the
relations defined by \(\psi\) and \(\hat\psi\) coincide for every value of
\(\bar z\), so there \(\varphi\) is equivalent to \(\varphi_1=Q^{\mathrm
H}x\,x'\,y\,y'\,\hat\psi\).  For an equivalence relation \(e\) on
\(\{1,\ldots,k\}\) let \(\delta_e(\bar z)\) be the conjunction of
\(z_i=z_j\) for \(i\mathrel ej\) and \(z_i\neq z_j\) otherwise (for
\(k=0\) there is one \(e\), and \(\delta_e\) is \(\exists x\,x=x\)); a
tuple \(\bar a\) has \emph{type} \(e\) if \(\delta_e(\bar a)\) holds.

Claim: if \(\kappa\) and \(\lambda\) are infinite, \(\bar a\in\kappa^k\)
and \(\bar b\in\lambda^k\) have the same type \(e\), and
\(\kappa\models\varphi_1[\bar a]\), then \(\lambda\models\varphi_1[\bar
b]\).  Let \(f,g:\kappa\to\kappa\) witness \(\varphi_1[\bar a]\), and
consider the structure \(\mathfrak K=(\kappa,f,g,a_1,\ldots,a_k)\) for
new symbols \(F,G,c_1,\ldots,c_k\).  It satisfies the first-order sentence
\[
 \gamma=\forall x\,\forall x'\,\hat\psi(x,x',F(x),G(x'),\bar c)\wedge
 \delta_e(\bar c).
\]
\(\mathfrak K\) is infinite and its vocabulary is finite, so by
\cref{thm:upward-ls} there is \(\mathfrak L'\equiv\mathfrak K\) of
cardinality \(|\lambda|\).  \(\mathfrak L'\models\gamma\), so with
\(\bar b'=\bar c^{\mathfrak L'}\) of type \(e\) the functions
\(F^{\mathfrak L'},G^{\mathfrak L'}\) witness \(\varphi_1[\bar b']\) in the
universe \(L'\) of \(\mathfrak L'\).  The sets \(L'-\{b'_1,\ldots,b'_k\}\)
and \(\lambda-\{b_1,\ldots,b_k\}\) are infinite of cardinality
\(|\lambda|\), and \(\bar b'\), \(\bar b\) have the same type, so there is
a bijection \(\beta:L'\to\lambda\) with \(\beta(b'_i)=b_i\) for all \(i\).
A bijection preserves the truth of the first-order pure identity formula
\(\hat\psi\), so \(\beta F^{\mathfrak L'}\beta^{-1}\) and \(\beta
G^{\mathfrak L'}\beta^{-1}\) witness \(\varphi_1[\bar b]\) in \(\lambda\).

Let \(E\) be the set of types \(e\) for which \(\kappa\models\varphi_1[\bar
a]\) for some infinite \(\kappa\) and some \(\bar a\) of type \(e\), and
\(\hat\varphi=\bigvee_{e\in E}\delta_e\) (the empty disjunction being
\(\neg\exists x\,x=x\)).  By the claim, for every infinite \(\kappa\) and
every \(\bar a\), \(\kappa\models\varphi[\bar a]\) iff
\(\kappa\models\varphi_1[\bar a]\) iff the type of \(\bar a\) is in \(E\)
iff \(\kappa\models\hat\varphi[\bar a]\).  The axiom of choice enters
through \cref{thm:upward-ls}.
\end{proof}
\end{annexitem}

\begin{annexitem}{annex:p03-w03-p03-042-mundici-q-alpha}{thm:p03-w03-p03-042-mundici-q-alpha}{Which \(Q_\alpha\) the Henkin quantifier defines}
\begin{proof}
For \(\alpha=0\): translate \(\mathscr L(Q_0)\) into \(\mathscr
L(Q^{\mathrm H})\) by induction on formulas as in the proof of
\cref{thm:p03-w03-p03-042-mundici-walkoe}, replacing \(Q_0x\,\varphi^\ast\)
by the right side of \cref{thm:p03-w03-p03-042-mundici-q-zero} for
\(\varphi^\ast\); so \(\mathscr L(Q^{\mathrm H})\geq\mathscr L(Q_0)\).

Let \(\alpha>0\) and suppose \(\mathscr L(Q^{\mathrm H})\geq\mathscr
L(Q_\alpha)\).  The sentence \(Q_\alpha x\,(x=x)\) of the empty vocabulary
then has the same models as a sentence \(\varphi\) of \(\mathscr
L(Q^{\mathrm H})\) of the empty vocabulary, and by
\cref{thm:p03-w03-p03-042-mundici-pure-identity} there is a first-order
sentence \(\hat\varphi\) of the empty vocabulary with
\(\kappa\models Q_\alpha x\,(x=x)\leftrightarrow\hat\varphi\) for every
infinite \(\kappa\).  Take \(\kappa=\aleph_\alpha\).  By
\cref{thm:downward-ls} there is \(\mathfrak N\preccurlyeq\aleph_\alpha\) (as
a structure of the empty vocabulary) of cardinality at most \(\aleph_0\).
For each \(n\) the first-order sentence ``there are at least \(n\)
elements'' holds in \(\aleph_\alpha\), hence in \(\mathfrak N\), so
\(\mathfrak N\) is infinite and has cardinality \(\aleph_0\); a bijection
of \(\mathfrak N\) onto \(\omega\) is an isomorphism of structures of the
empty vocabulary.  Hence \(\omega\equiv\aleph_\alpha\), and
\(\aleph_\alpha\models\hat\varphi\) (since
\(\aleph_\alpha\models Q_\alpha x\,(x=x)\)) gives \(\omega\models
\hat\varphi\), so \(\omega\models Q_\alpha x\,(x=x)\), that is,
\(\aleph_0\geq\aleph_\alpha\), against \(\alpha>0\).
\end{proof}
\end{annexitem}

\begin{annexitem}{annex:p03-w03-p03-042-kq-effective}{thm:p03-w03-p03-042-kq-effective}{Effectivity for \(L(Q)\)}
\begin{proof}
All procedures below query the oracle \((\Sigma_L,\kappa_L)\) finitely
often and halt.

(1) \emph{Formulas and sentences.}  Given a code, decode the finite
sequence of symbol numbers; for each, ask whether it is a logical symbol,
a variable, or in \(\Sigma_L\) and of which kind and arity.  Terms and
formulas are recognized by the formation rules of
\cref{def:p02-w02c-p02-003-kq-language}: a string is a term iff it is a
variable, a constant, or \(f\) of arity \(k\) followed by \(k\) terms,
which is checked by trying the finitely many ways to split the remainder;
formulas are checked likewise.  Free occurrences are computed by the
recursive clauses, so sentences are recognized.

\emph{Axioms.}  (0a): compute the prime components of \(\varphi\) (its
maximal subformulas that are not negations or conjunctions), assign them
sentence letters, and evaluate the resulting propositional formula under
all \(2^k\) assignments.  (0b): \(\varphi\) must have the form
\((\forall x)\psi\to\chi\).  If \(x\) is not free in \(\psi\), the test is
\(\chi=\psi\).  Otherwise the first free occurrence of \(x\) in \(\psi\)
corresponds to a position in \(\chi\) where a term \(t\) begins; \(t\) is
determined by that position, and the test is that \(t\) is substitutable
for \(x\) in \(\psi\) and \(\chi=\psi(t/x)\).  (0c) is checked in the same
way.  (0d), (0e): match the displayed shape and test that \(x\) is not free
in \(\psi\).  (0f), (0g): match the shape; for (0g) compare \(\alpha\) and
\(\alpha'\) symbol by symbol, allowing \(x\) to become \(y\) at variable
positions.  Axioms~1, 2 and 4 are shape matches with the variable
distinctness conditions.  Axiom~3: match
\((Qx)\varphi\leftrightarrow(Qy)\varphi'\) (with \(\leftrightarrow\)
unfolded into \(\neg,\land\)), test that \(y\) does not occur in
\(\varphi\), and test \(\varphi'=\varphi(y/x)\).  Each test inspects a
finite string, so the set of axioms is recursive in the oracle.

(2) Given a finite sequence \(\varphi_1,\ldots,\varphi_n\), test each
\(\varphi_i\): an axiom, or some \(j,k<i\) with \(\varphi_k=(\varphi_j\to
\varphi_i)\) (modus ponens), or some \(j<i\) and variable \(x\) with
\(\varphi_i=(\forall x)\varphi_j\) (generalization).  Finitely many cases.

(3) List all pairs \((p,\varphi)\) of codes and output \(\varphi\) when
\(p\) is a deduction of \(\varphi\) from the empty set; this enumerates the
provable formulas, so they are recursively enumerable in the oracle.
Intersecting with the recursive set of sentences gives the provable
sentences.  By \cref{thm:p02-w02c-p02-003-keisler-completeness} (Keisler's
Corollary~2.10, first half), a sentence is provable iff it is valid, so
the valid sentences are recursively enumerable in the oracle.  If the
oracle is recursive, relative recursive enumerability is recursive
enumerability.  If \(L\) is finite, its presentation can be chosen with
finite \(\Sigma_L\) and finite \(\kappa_L\), both recursive.
\end{proof}
\end{annexitem}

\begin{annexitem}{annex:p03-w03-p03-042-bridge}{thm:p03-w03-p03-042-bridge}{Correspondence between the two presentations}
\begin{proof}
(1) If \((M,A)\in\mathsf K(Q)\) and \(\varphi:M\to M'\) is a bijection,
the invariance clause of unlimited quantifiers gives
\(Q_{M'}(\varphi[A])=Q_M(A)=\vee\), so \((M',\varphi[A])\in\mathsf K(Q)\);
an isomorphism of pairs \((M,A)\cong(M',A')\) is such a bijection with
\(\varphi[A]=A'\).  Conversely, if \(\mathcal K\) is isomorphism-closed,
\(Q^{\mathcal K}_M\) is invariant under permutations of \(M\) and
\(Q^{\mathcal K}_M(A)=Q^{\mathcal K}_{M'}(\varphi[A])\) for bijections
\(\varphi\), so it is an unlimited quantifier on nonempty sets.  The two
maps are mutually inverse by their definitions.

(2) Induction on \(\mathrm Z\).  Atomic \(\mathrm F_j(\mathrm x_{i_1},
\ldots,\mathrm x_{i_k})\): \(\mathit{val}_{MI}\) is
\([\mathrm F_j]_M\) at the padded sequence, which is \(\vee\) iff the tuple
of values lies in \(R_j^{\mathcal M}\).  Identity: both sides compare
values.  Stroke: \(\mathit{val}(\mathrm Z_1|\mathrm Z_2)=\vee\) iff not
both are \(\vee\), iff \(\mathcal M,s\entails\neg(\mathrm Z_1^\circ\land
\mathrm Z_2^\circ)\).  Quantifier: by induction the set in clause~4 of
\cref{def:p03-w03-p03-042-calculus-s} equals \(\{a\in I:\mathcal
M,s[\mathrm x_i\mapsto a]\entails\mathrm Z_1^\circ\}\), and the clause of
\cref{def:p02-generalized-quantifier} for \(\mathsf K(Q^j)\) asks exactly
whether \((I,\text{that set})\in\mathsf K(Q^j)\), i.e.\ whether
\(Q^j_I\) of it is \(\vee\).

(3) \(({\sim}P)_I(A)=\vee\) iff \(P_I(A)=\wedge\) iff \(\overline A>
\aleph_0\) (cardinals are comparable by the axiom of choice), which is
the clause of \cref{def:p03-q-one}; so \(Q_1=\mathsf K({\sim}P)\) and the
determining function is \(T(\mathfrak m,\mathfrak n)=\vee\) iff
\(\mathfrak m>\aleph_0\).  On a denumerable \(I\) every \(A\) is
countable, so the value is always \(\wedge\); hence the set of \(n\) with
\(T(n,\aleph_0)=\vee\) is empty and \(Q_1\) lacks property (E).  If
\(\overline I>\aleph_0\), take \(A\) of power \(\aleph_0\) with
\(\overline{I-A}=\overline I\) and \(B\) of power \(\aleph_1\) with
\(\overline{I-B}=\overline I\) (split \(I\) into two parts of power
\(\overline I\) and choose inside one of them).  Every
\(\overset{(m)}\Sigma_I\) and \(\overset{(n)}\Pi_I\) takes the value
\(\wedge\) at both \(A\) and \(B\), since all four cardinals involved are
infinite, so every Boolean polynomial in them takes equal values at \(A\)
and \(B\); but \(Q_1\) gives \(\wedge\) at \(A\) and \(\vee\) at \(B\).
So \(Q_1\) limited to \(I\) is not numerical.
\end{proof}
\end{annexitem}

\begin{annexitem}{annex:p03-w03-p03-042-q-one-host-form}{thm:p03-w03-p03-042-q-one-host-form}{The host statement on \(L(Q_1)\)}
\emph{Reduct lemma.}  If \(L_0\subseteq L\), \(\mathfrak A\) is an
\(L\)-structure and \(\varphi\) a formula of \(L_0(Q)\), then
\(\mathfrak A\models\varphi[s]\) iff \(\mathfrak A{\restriction}L_0\models
\varphi[s]\) in standard semantics.  Induction on \(\varphi\): atomic
formulas use only symbols of \(L_0\), which are interpreted alike; the
connective clauses are immediate; for \(\exists,\forall\) and \(Q\) the
two structures have the same universe, so by induction the defining sets
coincide, and uncountability of a subset of the common universe does not
depend on the language.
\begin{proof}
By \cref{def:p03-w03-p03-042-keisler-bridge}, standard models of
\(L(Q)\) over \(\tau\) are the \(\tau\)-structures with the \(Q_1\) clause,
so \(\tau\)-validity and satisfiability in the host sense are Keisler's
notions for the language \(\tau\).

(1) Let \(\sigma\) be derivable over \(\tau\).  A derivation uses finitely
many symbols; let \(L_1\) be the finite language of those symbols and of
\(\sigma\).  Every line is a formula of \(L_1(Q)\), every axiom used is an
axiom of \(L_1(Q)\) (the schemes and their side conditions refer only to
the formulas involved), and the rules are the same, so the derivation is
one over \(L_1\).  By part~4 of \cref{thm:p02-w02c-p02-003-kq-soundness}
for the countable \(L_1\), \(\sigma\) holds in every standard
\(L_1\)-model, hence by the reduct lemma in every \(\tau\)-structure.
Conversely let \(\sigma\) be \(\tau\)-valid and \(L_0\) the finite set of
symbols of \(\sigma\).  Every \(L_0\)-structure expands to a
\(\tau\)-structure (interpret the other symbols arbitrarily on the
nonempty universe), in which \(\sigma\) holds, so by the reduct lemma
\(\sigma\) holds in every \(L_0\)-structure.  By
\cref{thm:p02-w02c-p02-003-keisler-completeness} for \(L_0\), \(\sigma\)
is derivable over \(L_0\), and this derivation is one over \(\tau\).

(2) Let \(L_0\) be the set of symbols occurring in the countable set
\(\Sigma\); it is countable, since fixing an enumeration of \(\Sigma\) and listing
the symbols of each sentence in order of occurrence lists \(L_0\).  If every finite \(\Sigma_0\subseteq\Sigma\)
has a \(\tau\)-model, its reduct to \(L_0\) is a standard \(L_0\)-model
of \(\Sigma_0\).  By \cref{thm:p02-w02c-p02-003-keisler-compactness}
\(\Sigma\) has a standard \(L_0\)-model, which expands to a
\(\tau\)-structure; by the reduct lemma it is a model of \(\Sigma\).

(3) Let \(\tau\) be countable and \(\Sigma\) satisfiable.  By part~4 of
\cref{thm:p02-w02c-p02-003-kq-soundness}, \(\Sigma\) is consistent, so by
\cref{thm:p02-w02c-p02-003-keisler-completeness} it has a model of power
at most \(\aleph_1\).  If \(\Sigma\cup\{Q_1x\,(x=x)\}\) is satisfiable it is
consistent, and the same theorem gives a model of \(\Sigma\cup
\{Q_1x\,(x=x)\}\) of power exactly \(\aleph_1\).  Conversely a model of
power \(\aleph_1\) has an uncountable universe, so it satisfies
\(Q_1x\,(x=x)\).

(4) With a unary \(P\) and distinct constants \(c_\alpha\)
(\(\alpha<\omega_1\)) in \(\tau\), the set \(\Delta\) of
\cref{def:p02-w02c-p02-003-kq-example} has power \(\aleph_1\), every finite
subset has a model, and \(\Delta\) has none; every finite subset is
consistent by soundness, so \(\Delta\) is consistent.
\end{proof}
\end{annexitem}

\begin{annexitem}{annex:p03-w03-p03-042-monadic-bridge}{thm:p03-w03-p03-042-monadic-bridge}{Theorems~10 and~11 in the handbook's notation}
\begin{proof}
\emph{Translation.}  Given unary generalized quantifiers
\(\mathcal K^2,\ldots,\mathcal K^s\), let \(Q^1=\exists\) and
\(Q^j=Q^{\mathcal K^j}\) (\(j\geq2\)) as in
\cref{def:p03-w03-p03-042-mostowski-bridge}, defined on nonempty sets and
on the empty set by any fixed value.  Map an equality-free formula
\(\varphi\) over \(\tau\) to a formula \(\mathrm Z_\varphi\) of
\((\mathrm S^M)\): \(R_h(v_i)\mapsto\mathrm F_h(\mathrm x_i)\);
\(\neg\psi\mapsto\mathrm Z_\psi|\mathrm Z_\psi\);
\(\psi\land\chi\mapsto(\mathrm Z_\psi|\mathrm Z_\chi)|(\mathrm Z_\psi|
\mathrm Z_\chi)\), and the other connectives through these;
\(\exists v_i\,\psi\mapsto(\mathbf Q^1\mathrm x_i)\mathrm Z_\psi\);
\(\forall v_i\,\psi\mapsto{\sim}(\mathbf Q^1\mathrm x_i){\sim}\mathrm
Z_\psi\); and the quantifier symbol of \(\mathcal K^j\) binding \(v_i\)
\(\mapsto\mathbf Q^j\mathrm x_i\).  The map is computable.  For a nonempty
\(I\) and an \(I\)-valuation \(M\), let \(\mathcal M\) and \(s\) be as in
\cref{thm:p03-w03-p03-042-bridge}(2).  Induction on \(\varphi\) gives
\(\mathit{val}_{MI}(\mathrm Z_\varphi)=\vee\) iff \(\mathcal M,s\entails
\varphi\): atoms and connectives by the clauses; \(\exists\) because
\(Q^1_I(X)=\vee\) iff \(X\neq\varnothing\); \(\forall\) as the dual; and
the generalized quantifier clause because \(Q^{\mathcal K^j}_I(X)=\vee\)
iff \((I,X)\in\mathcal K^j\).  Every \(\tau\)-structure with an
assignment arises from such \(I,M\): take its universe for \(I\), its
relations for the \([\mathrm F_h]_M\) and its assignment for the
\([\mathrm x_i]_M\).  As every formula is true in the empty set,
\(\varphi\) is valid iff \(\mathrm Z_\varphi\) is true.

(1) With \(\mathcal K^j=\mathsf K_{\geq\mathfrak m_j}\) and \(\mathfrak
m_1=1\), the quantifiers \(Q^1,\ldots,Q^s\) are those of
\cref{thm:p03-w03-p03-042-theorem-ten}: \(Q^j_I(X)=\vee\) iff
\(\overline X\geq\mathfrak m_j\), and the hypotheses on the cardinals
hold.  So validity of \(\varphi\) is decided by computing \(\mathrm
Z_\varphi\) and applying the procedure of
\cref{annex:p03-w03-p03-042-theorem-ten}.  For the particular case take
\(\mathfrak m_2=\aleph_0\) and \(\mathfrak m_3=\aleph_1\): ``infinitely
many'' is ``at least \(\aleph_0\)'', and ``uncountably many'' is ``at
least \(\aleph_1\)'', because by the axiom of choice a power that is not
at most \(\aleph_0\) is at least \(\aleph_1\), the least cardinal above
\(\aleph_0\).

(2) Apply the translation to the quantifiers \(\exists,\mathcal
K^1,\ldots,\mathcal K^s\), numbered \(1,\ldots,s+1\) in this order.  \(\mathrm
Z_\varphi\) is a formula of \((\mathrm S^M_n)\), and
\cref{thm:p03-w03-p03-042-theorem-eleven} makes the set of true formulas
of \((\mathrm S^M_n)\) recursive; its preimage under the computable map
\(\varphi\mapsto\mathrm Z_\varphi\) is the set of valid formulas.
\end{proof}
\end{annexitem}
